\chapter{How Idiomatic Rust Is Actually Written: Architecture, Patterns, and Anti-Patterns}
\label{ch08:idiomatic-rust}

Chapters~3 through~7 were about single functions: what each byte on the stack and heap is doing, how borrows and lifetimes are checked, how enums, traits, and iterators compile. This chapter is about the next problem: putting a few thousand lines together.

When you build a five-thousand-line service of your own, a different set of errors appears. You try to wire up a hierarchical chart of accounts where each child account points back to its parent and each parent holds a list of children, and the borrow checker blocks every edge. You reach for \texttt{Rc<RefCell<Account>>} to make the compiler stop complaining, only to watch a subtle re-entrant call panic at runtime. You write a \texttt{SettlementEngine} struct with ten fields, call \texttt{self.\allowbreak{}apply\_posting(\&self.fee\_schedule)} inside a \texttt{\&mut self} method, and hit compiler error \texttt{E0502}: cannot borrow \texttt{self.fee\_schedule} as immutable because it is also borrowed as mutable. You pass \texttt{\&String} and \texttt{\&Vec<Posting>} across helper functions, sprinkle defensive \texttt{if amount <= 0} checks at the top of every method, and end up with more checks than logic.

\begin{intuitionbox}
Fighting the borrow checker is almost never a sign that you forgot a lifetime annotation. Nine times out of ten, it is an architectural diagnostic: you are sketching a \emph{Java object graph} --- cyclic pointers, shared mutable state, runtime boolean flags, and unvalidated primitive arguments --- and asking Rust to certify that no two pointers alias mutably. Idiomatic Rust does not patch Java architectures with extra lifetime syntax or \texttt{RefCell} wrappers. It reshapes data ownership, state transitions, and module boundaries so that the hardware layout is flat and invalid states cannot even be written down.
\end{intuitionbox}

This chapter answers the question at the heart of Part~III: \emph{how do experienced systems engineers actually see a problem through the Rust lens, and how should production Rust be structured?} Every section contrasts a familiar instinct from Java~6 or C++ against the idiomatic Rust design that replaces it, anchored in the financial ledger, payment routing, and settlement domain we will take all the way to a live Axum, SQLx, and PostgreSQL backend in Part~V.

\chapterline

% ======================================================================
\section{Curing ``Java-in-Rust Syndrome'': Data-Oriented Arenas and Typed Index Handles}
\label{ch08:arenas-and-handles}
% ======================================================================

Consider a core piece of any corporate treasury or double-entry accounting system: a \emph{hierarchical chart of accounts}. A multinational corporation holds a root treasury account, regional sub-accounts under that root, and operating accounts under each region. Whenever an operating account settles a wire transfer, two things happen: the sender and receiver operating balances change, and each ancestor account along the tree updates its rolled-up exposure total so risk limits can be checked in $O(1)$ time.

In Java~6, every developer draws and codes this the exact same way without thinking twice:

\begin{lstlisting}[language=Java]
// Java 6: bidirectional heap pointers everywhere
public class AccountNode {
    private final String code;
    private long balanceCents;
    private long subtreeExposureCents;
    private AccountNode parent; // null for root
    private final List<AccountNode> children =
        new ArrayList<AccountNode>();

    public void addChild(AccountNode child) {
        child.parent = this;
        this.children.add(child);
    }
}
\end{lstlisting}

In Java, object identity and object access are fused into a single concept: a garbage-collected heap pointer. Every \texttt{AccountNode} floats independently on the heap, \texttt{parent} points upward, \texttt{children} points downward, an order-book map holds another reference from the side, and an audit listener holds a fourth. Because the JVM makes no compile-time promises about aliasing, any method can mutate any reachable node at any instant.

Here is what happens when an engineer translates \texttt{AccountNode} line for line into Rust.

\subsection{Why Direct Translation Produces \texttt{Rc<RefCell<T>>} Soup}

The first attempt --- storing \texttt{parent:\allowbreak\ Option<\&'a mut AccountNode<'a>>} alongside \texttt{children:\allowbreak\ Vec<AccountNode<'a>>} --- is rejected immediately by the compiler. If the parent owns the child inside \texttt{children}, and the child simultaneously holds a mutable reference \texttt{\&mut} back to its parent, you have two active paths capable of mutating the parent at the same time: the owner path from above and the \texttt{parent} pointer from below. That directly violates Chapter~3's \emph{Aliasing XOR Mutability} law. In addition, moving the parent \texttt{Vec} during reallocation would invalidate every child's raw address.

Frustrated, the newcomer discovers Chapter~3's escape hatches for shared ownership and interior mutability --- \texttt{Rc<T>}, \texttt{Weak<T>}, and \texttt{RefCell<T>} --- and wraps every node in sight:

\begin{lstlisting}[language=Rust]
use std::cell::RefCell;
use std::rc::{Rc, Weak};

// Anti-pattern: recreating a Java object graph via Rc<RefCell<...>>
type SharedAccount = Rc<RefCell<BadAccountNode>>;

struct BadAccountNode {
    code: String,
    balance_cents: i64,
    subtree_exposure_cents: i64,
    parent: Option<Weak<RefCell<BadAccountNode>>>,
    children: Vec<SharedAccount>,
}
\end{lstlisting}

It compiles. But wrapping a domain graph in \texttt{Rc<RefCell<T>>} forfeits almost every advantage you came to Rust to get, while introducing failure modes Java~6 never had:

\begin{enumerate}
  \item \textbf{Heap fragmentation and CPU cache misses.} Every \texttt{Rc::new(RefCell::new(...))} performs an individual allocator request. Inside each heap block sits a 16-byte reference-count header (\texttt{strong: usize}, \texttt{weak: usize}) plus an 8-byte \texttt{BorrowFlag} (\texttt{isize}) before your fields even begin. Walking 10,000 accounts chases arbitrary 64-bit addresses scattered across megabytes of virtual memory, stalling the CPU pipeline for $\sim 100\text{--}200$ cycles on L3/DRAM misses instead of streaming sequential 64-byte cache lines.
  \item \textbf{Accidental reference-cycle memory leaks.} If a developer later adds a cross-reference field \texttt{counterparty:\allowbreak\ Option<SharedAccount>} using strong \texttt{Rc} instead of \texttt{Weak}, two accounts pointing at each other keep their strong counts $\ge 1$ forever. Rust has no tracing garbage collector; those nodes and everything they own leak silently for the lifetime of the process.
  \item \textbf{Thread hostility.} \texttt{Rc} and \texttt{RefCell} use non-atomic integer increments, so the compiler strips \texttt{Send} and \texttt{Sync} from the entire graph. You cannot pass this chart of accounts into a multi-threaded Tokio HTTP server in Part~V without rewriting every node to \texttt{Arc<Mutex<T>>} or \texttt{Arc<RwLock<T>>}, turning every field read into an atomic lock acquisition and inviting deadlocks when two transfers lock nodes in opposite orders.
  \item \textbf{Delayed runtime \texttt{BorrowMutError} panics.} \texttt{RefCell} does not remove the aliasing rule --- it merely moves enforcement from compile time to runtime.
\end{enumerate}

\begin{trapbox}
Here is how \texttt{Rc<RefCell<T>>} panics at runtime in an ordinary helper call:
\begin{lstlisting}[language=Rust]
fn debit_and_recompute(acct: &SharedAccount, cents: i64) {
    // 1. Acquire exclusive runtime borrow (borrow_flag becomes -1)
    let mut node = acct.borrow_mut();
    node.balance_cents -= cents;

    // 2. Ask parent to recompute exposure across all its children
    if let Some(ref weak_parent) = node.parent {
        if let Some(parent_rc) = weak_parent.upgrade() {
            recompute_parent(&parent_rc); // <-- CRASH HIDING HERE!
        }
    }
}

fn recompute_parent(parent: &SharedAccount) {
    let mut p = parent.borrow_mut();
    p.subtree_exposure_cents = p.children
        .iter()
        // One of p.children IS `acct`, still mutably borrowed above!
        // child.borrow() panics: "RefCell already mutably borrowed"
        .map(|child| child.borrow().balance_cents.abs())
        .sum();
}
\end{lstlisting}
The compiler approves both functions without a single warning. Everything passes unit tests on root accounts that have \texttt{parent: None}. Then, in production, a transfer hits a child account: \texttt{debit\_and\_recompute} holds \texttt{node = acct.borrow\_mut()} on the stack while \texttt{recompute\_parent} iterates over \texttt{p.children} and calls \texttt{child.borrow()} on that exact same child. The thread immediately panics with \texttt{BorrowMutError}.
\end{trapbox}

\subsection{The Rust Architectural Shift: Separate Identity from Ownership}

Ask what \texttt{parent} and \texttt{children} in our chart of accounts actually mean. Does a child account \emph{own the lifetime of} its parent account's memory? No. In a database, the \texttt{accounts} table owns all rows, and \texttt{parent\_account\_id} is simply an integer foreign key naming another row in that same table.

Idiomatic Rust applies that exact relational insight directly to in-memory data structures:

\begin{focusbox}
\textbf{Separate Identity from Ownership.} Whenever entities form graphs, trees with parent pointers, doubly linked queues, or cross-referenced registries, give \textbf{one container} (an \emph{arena}, backed by a contiguous \texttt{Vec}) exclusive ownership of all nodes, and represent edges between nodes as small, \texttt{Copy}, strongly typed \textbf{index handles} rather than pointers or reference-counted wrappers.
\end{focusbox}

Every problem from the \texttt{Rc<RefCell<T>>} list goes away with this shift:
\begin{itemize}
  \item \textbf{One clear owner, automatic cleanup.} When the \texttt{AccountArena} struct goes out of scope, dropping its single \texttt{Vec} frees every node at once. Cycles among integer handles are just numbers; they can never leak memory.
  \item \textbf{Hardware sympathy.} Every \texttt{AccountNode} sits back to back in one contiguous slab of RAM. Scanning balances streams straight through the CPU hardware prefetcher, and handles shrink from 8-byte aligned pointers (plus 24 bytes of \texttt{Rc}/\texttt{RefCell} metadata) down to compact 4-byte or 8-byte values.
  \item \textbf{Automatic \texttt{Send} and \texttt{Sync}.} Plain structs holding integers, \texttt{Vec}s, and \texttt{String}s are automatically \texttt{Send + Sync} with zero atomic refcount overhead.
  \item \textbf{Compile-time borrow checking restored.} Instead of hidden per-node \texttt{RefCell} flags, borrowing happens at the arena boundary: \texttt{arena.get(key)} takes \texttt{\&self} (allowing unlimited simultaneous readers), and \texttt{arena.get\_mut(key)} takes \texttt{\&mut self} (making mutation explicit in the function signature).
\end{itemize}

\begin{figure}[H]
\centering
\begin{tikzpicture}[
    font=\small,
    box/.style={draw, rounded corners=2pt, thick, align=left},
    ptr/.style={-{Stealth[length=2.3mm,width=1.7mm]}, thick}
]
  % Left panel: Java / Rc<RefCell<T>> scattered heap
  \draw[draw=trapred!70, fill=trapred!4, rounded corners=4pt, thick] (0,0) rectangle (6.45,5.6);
  \node[font=\bfseries\small, text=trapred!85!black] at (3.22,5.25) {Java / \texttt{Rc<RefCell<Node>>} Graph};
  \node[font=\scriptsize, text=gray!80!black] at (3.22,4.88) {Scattered heap allocations + refcount/borrow headers};

  \node[box, draw=trapred!80, fill=white, font=\scriptsize\ttfamily, inner sep=3.5pt] (rc0) at (1.75,3.65) {
    strong:2 | weak:2\\
    borrow\_flag: -1\\
    Account "ROOT"
  };
  \node[box, draw=trapred!80, fill=white, font=\scriptsize\ttfamily, inner sep=3.5pt] (rc1) at (4.75,2.45) {
    strong:1 | weak:0\\
    borrow\_flag: 0\\
    Account "EU-OPS"
  };
  \node[box, draw=trapred!80, fill=white, font=\scriptsize\ttfamily, inner sep=3.5pt] (rc2) at (1.95,1.15) {
    strong:1 | weak:0\\
    borrow\_flag: 0\\
    Account "US-OPS"
  };

  \draw[ptr, trapred!85!black, bend left=16] (rc0.east) to node[above, sloped, font=\tiny] {\texttt{Rc} ptr} (rc1.north);
  \draw[ptr, trapred!65!black, dashed, bend left=16] (rc1.west) to node[below, sloped, font=\tiny] {\texttt{Weak}} (rc0.south east);
  \draw[ptr, trapred!85!black, bend right=14] (rc0.south) to node[left, font=\tiny] {\texttt{Rc}} (rc2.north);
  \draw[ptr, trapred!65!black, dashed, bend right=18] (rc2.east) to node[right, font=\tiny] {\texttt{Weak}} (rc0.south);
  \node[font=\scriptsize, text=trapred!85!black] at (3.22,0.32) {Cache misses + cycle leaks + runtime \texttt{BorrowMutError}};

  % Right panel: Flat Contiguous Arena + Typed Handles
  \draw[draw=targetgreen!80!black, fill=targetgreen!4, rounded corners=4pt, thick] (6.85,0) rectangle (13.8,5.6);
  \node[font=\bfseries\small, text=targetgreen!70!black] at (10.32,5.25) {Idiomatic Rust Arena + \texttt{AccountKey}};
  \node[font=\scriptsize, text=gray!80!black] at (10.32,4.88) {Single contiguous \texttt{Vec<Slot>} slab in L1/L2 cache};

  % Contiguous array slots
  \draw[draw=sourceblue!85!black, fill=white, thick] (7.2,3.45) rectangle (13.45,4.45);
  \node[font=\scriptsize\ttfamily, anchor=west] at (7.3,4.15) {[0] gen:1 | "ROOT"   | bal: 9000};
  \node[font=\scriptsize\ttfamily, text=sourceblue!90!black, anchor=west] at (7.3,3.72) {    parent: None, children: [(1,1), (2,1)]};

  \draw[draw=sourceblue!85!black, fill=white, thick] (7.2,2.45) rectangle (13.45,3.45);
  \node[font=\scriptsize\ttfamily, anchor=west] at (7.3,3.15) {[1] gen:1 | "EU-OPS" | bal: 4000};
  \node[font=\scriptsize\ttfamily, text=sourceblue!90!black, anchor=west] at (7.3,2.72) {    parent: Some(Key(0,1)), children: []};

  \draw[draw=sourceblue!85!black, fill=white, thick] (7.2,1.45) rectangle (13.45,2.45);
  \node[font=\scriptsize\ttfamily, anchor=west] at (7.3,2.15) {[2] gen:1 | "US-OPS" | bal: 5000};
  \node[font=\scriptsize\ttfamily, text=sourceblue!90!black, anchor=west] at (7.3,1.72) {    parent: Some(Key(0,1)), children: []};

  % Handle badge at bottom
  \draw[draw=structurepurple, fill=structurepurple!10, rounded corners=2pt, thick] (7.55,0.62) rectangle (13.1,1.18);
  \node[font=\scriptsize\ttfamily, text=structurepurple!85!black] at (10.32,0.90) {AccountKey \{ index: u32, generation: u32 \} (8 B Copy)};
  \node[font=\scriptsize, text=targetgreen!75!black] at (10.32,0.30) {Sequential RAM + 0 refcounts + stale-handle ABA check};
\end{tikzpicture}
\caption{Pointer-based object graphs versus a flat generational arena. Left: \texttt{Rc<RefCell<T>>} scatters nodes across separate heap allocations with 24 bytes of refcount and borrow-flag overhead per node. Right: one \texttt{Vec} owns all nodes contiguously while lightweight \texttt{Copy} keys encode relationships as data.}
\label{ch08:fig-arena-vs-rc}
\end{figure}

\subsection{Generational Indices: Solving the Stale-Handle (ABA) Problem}

A sharp C programmer looking at Figure~\ref{ch08:fig-arena-vs-rc} immediately asks an important question: \emph{``Wait --- if an index handle is just \texttt{struct AccountKey(u32)}, what happens when account \texttt{1} is closed, removed from the arena, and its slot is later reused for a brand-new customer?''}

If a queued settlement batch still holds the old \texttt{AccountKey(1)}, indexing into slot \texttt{1} will not corrupt raw memory the way a C dangling pointer does, yet it will commit an equally nasty \emph{logical} bug: crediting money to the wrong customer. In systems programming, this is the classic \textbf{ABA problem}: slot $1$ went from Customer $A$, to free, to Customer $B$, and a naive index cannot tell the second occupant apart from the first.

The cure --- used for entity handles in Rust game engines, in compilers (\texttt{rustc} itself stores AST, HIR, and MIR graphs in arenas), and the widely used \texttt{slotmap} crate --- is a \textbf{generational index}:

\begin{definition}[Generational Handle and Slot Map]
A \emph{generational arena} pairs each array slot $i \in \{0, \dots, n-1\}$ with a monotonically increasing counter $\text{gen}_i \in \N$. Every handle handed to callers is an 8-byte pair $(\text{index}: \texttt{u32},\; \text{generation}: \texttt{u32})$.
\begin{itemize}
  \item On \textbf{lookup} of $(\text{idx}, g)$, the arena checks both $\text{idx} < n$ and that slot $\text{idx}$ is currently \texttt{Occupied} with $\text{gen}_{\text{idx}} == g$. If either fails, lookup cleanly returns \texttt{None}.
  \item On \textbf{removal} of $(\text{idx}, g)$, slot $\text{idx}$ transitions to \texttt{Vacant}, increments its generation to $g + 1$ (using wrapping or checked arithmetic), and pushes $\text{idx}$ onto an embedded singly linked free-list of reusable slots.
\end{itemize}
Any stale handle still holding $(\text{idx}, g)$ will forever see $g \neq g + 1$ and fail closed as \texttt{None}, even after the slot is recycled millions of times.
\end{definition}

Notice how neatly the free-list fits inside the vacant slots themselves without a second \texttt{Vec}: when a slot is empty, its payload space stores \texttt{next\_free: Option<u32>} pointing to the next available index.

\subsection{Mutating Two Nodes Safely with \texttt{split\_at\_mut}}

A second practical question surfaces the moment you write \texttt{transfer(sender\_key, receiver\_key, cents)} on an arena: \emph{how do you obtain \texttt{\&mut AccountNode} to both the sender and receiver at the same time?}

If you write two consecutive calls --- \texttt{let s = arena.get\_mut(sender)?;} followed by \texttt{let r = arena.get\_mut(receiver)?;} --- the compiler rejects the second line because \texttt{get\_mut(\&mut self, ...)} borrows the entire arena exclusively for the lifetime of \texttt{s}. And rightly so. What if a caller passed \texttt{sender == receiver}? Returning two \texttt{\&mut AccountNode} references pointing to the exact same slot would create aliasing mutable references --- instant undefined behavior in Rust's memory model.

How, then, do we give callers simultaneous mutable access to two distinct slots without \texttt{unsafe} code? The standard library provides \texttt{slice::split\_at\_mut(mid)}, which splits one mutable slice \texttt{\&mut [T]} at index \texttt{mid} into two provably non-overlapping sub-slices \texttt{(\&mut [T], \&mut [T])}: indices \texttt{0..mid} on the left, and indices \texttt{mid..len} on the right.

Let us put the entire generational \texttt{Arena<T>}, including \texttt{get\_two\_mut} and our hierarchical \texttt{Chart\allowbreak{}Of\allowbreak{}Accounts} ledger with ancestor exposure rollups, together in roughly 120 lines of 100-percent safe, production-grade Rust:

\begin{lstlisting}[language=Rust]
#[derive(Copy, Clone, Debug, PartialEq, Eq, Hash)]
pub struct AccountKey {
    index: u32,
    generation: u32,
}

enum Slot<T> {
    Occupied { generation: u32, value: T },
    Vacant { generation: u32, next_free: Option<u32> },
}

pub struct Arena<T> {
    slots: Vec<Slot<T>>,
    free_head: Option<u32>,
}

impl<T> Arena<T> {
    pub fn new() -> Self {
        Self { slots: Vec::new(), free_head: None }
    }

    pub fn insert(&mut self, value: T) -> AccountKey {
        if let Some(idx) = self.free_head {
            let slot = &mut self.slots[idx as usize];
            match *slot {
                Slot::Vacant { generation, next_free } => {
                    self.free_head = next_free;
                    *slot = Slot::Occupied { generation, value };
                    AccountKey { index: idx, generation }
                }
                Slot::Occupied { .. } => {
                    unreachable!("invariant: free_head names Vacant")
                }
            }
        } else {
            let idx = u32::try_from(self.slots.len())
                .expect("arena capacity exceeded u32::MAX slots");
            self.slots.push(Slot::Occupied { generation: 1, value });
            AccountKey { index: idx, generation: 1 }
        }
    }

    pub fn remove(&mut self, key: AccountKey) -> Option<T> {
        let slot = self.slots.get_mut(key.index as usize)?;
        match slot {
            Slot::Occupied { generation, .. }
                if *generation == key.generation =>
            {
                let next_gen = generation.saturating_add(1);
                let old = std::mem::replace(
                    slot,
                    Slot::Vacant {
                        generation: next_gen,
                        next_free: self.free_head,
                    },
                );
                self.free_head = Some(key.index);
                match old {
                    Slot::Occupied { value, .. } => Some(value),
                    Slot::Vacant { .. } => unreachable!(),
                }
            }
            _ => None,
        }
    }

    pub fn get(&self, key: AccountKey) -> Option<&T> {
        match self.slots.get(key.index as usize)? {
            Slot::Occupied { generation, value }
                if *generation == key.generation => Some(value),
            _ => None,
        }
    }

    pub fn get_mut(&mut self, key: AccountKey) -> Option<&mut T> {
        match self.slots.get_mut(key.index as usize)? {
            Slot::Occupied { generation, value }
                if *generation == key.generation => Some(value),
            _ => None,
        }
    }

    /// Safely borrow two distinct elements mutably at the same time.
    pub fn get_two_mut(
        &mut self,
        a: AccountKey,
        b: AccountKey,
    ) -> Option<(&mut T, &mut T)> {
        if a.index == b.index {
            return None; // Cannot mutably alias the same slot!
        }
        let (ai, bi) = (a.index as usize, b.index as usize);
        if ai >= self.slots.len() || bi >= self.slots.len() {
            return None;
        }
        let (slot_a, slot_b) = if ai < bi {
            let (left, right) = self.slots.split_at_mut(bi);
            (&mut left[ai], &mut right[0])
        } else {
            let (left, right) = self.slots.split_at_mut(ai);
            (&mut right[0], &mut left[bi])
        };
        match (slot_a, slot_b) {
            (
                Slot::Occupied { generation: ga, value: va },
                Slot::Occupied { generation: gb, value: vb },
            ) if *ga == a.generation && *gb == b.generation => {
                Some((va, vb))
            }
            _ => None,
        }
    }
}
\end{lstlisting}

Look at how clean the hierarchical chart of accounts becomes on top of \texttt{Arena<AccountNode>}:

\begin{lstlisting}[language=Rust]
pub struct AccountNode {
    pub code: String,
    pub balance_cents: i64,
    pub subtree_exposure_cents: i64,
    pub parent: Option<AccountKey>,
    pub children: Vec<AccountKey>,
}

pub struct ChartOfAccounts {
    nodes: Arena<AccountNode>,
}

impl ChartOfAccounts {
    pub fn transfer(
        &mut self,
        sender: AccountKey,
        receiver: AccountKey,
        cents: i64,
    ) -> Result<(), &'static str> {
        if cents <= 0 {
            return Err("transfer amount must be positive");
        }
        // Step 1: Mutate sender and receiver together via split_at_mut
        {
            let (s, r) = self
                .nodes
                .get_two_mut(sender, receiver)
                .ok_or("invalid or identical sender/receiver key")?;
            if s.balance_cents < cents {
                return Err("insufficient funds");
            }
            s.balance_cents -= cents;
            r.balance_cents += cents;
        } // <-- both &mut borrows of self.nodes end here!

        // Step 2: Walk upward from each touched leaf. Cannot panic.
        self.recompute_upward(Some(sender));
        self.recompute_upward(Some(receiver));
        Ok(())
    }

    // Recomputes exposure for `cursor` and then every ancestor above it.
    fn recompute_upward(&mut self, mut cursor: Option<AccountKey>) {
        while let Some(key) = cursor {
            let Some(node) = self.nodes.get(key) else { break };
            let next_parent = node.parent;
            // Sum own balance + children exposures via shared &self
            let own = node.balance_cents.abs();
            let child_sum: i64 = node
                .children
                .iter()
                .filter_map(|&ck| self.nodes.get(ck))
                .map(|c| c.subtree_exposure_cents)
                .sum();
            if let Some(node_mut) = self.nodes.get_mut(key) {
                node_mut.subtree_exposure_cents = own + child_sum;
            }
            cursor = next_parent;
        }
    }
}
\end{lstlisting}

\begin{whybox}
Trace why the re-entrancy crash from \texttt{Rc<RefCell<T>>} cannot happen here. In Step~1 of \texttt{transfer}, \texttt{s} and \texttt{r} borrow two slots of \texttt{self.nodes} only inside the inner block \texttt{\{ ... \}}, and the compiler ends both borrows at its closing brace. By Step~2 nothing is borrowed at all, so \texttt{recompute\_upward} is free to read any child slot through \texttt{get} and then write the current node through \texttt{get\_mut}, one step at a time. The ordering that \texttt{RefCell} checked at runtime (and got wrong in the trap above) is now checked by \texttt{rustc} before the program ever runs. Starting the walk at the touched leaf itself, rather than at its parent, also matters: the leaf's own \texttt{subtree\_exposure\_cents} must be refreshed before its parent sums it.
\end{whybox}

% ======================================================================
\section{Parse, Don't Validate: Zero-Cost Newtypes and Smart Constructors}
\label{ch08:parse-dont-validate}
% ======================================================================

Look back at the signature of \texttt{ChartOfAccounts::transfer} above. Notice how we already replaced raw \texttt{u32} account indices with \texttt{AccountKey}, yet the final parameter is still \texttt{cents: i64}, forcing the method to open with \texttt{if cents <= 0 \{ return Err(...); \}}. Now imagine scaling that style to an international wire-submission function written in the ``primitive obsession'' style common in Java~6 and C codebases:

\begin{lstlisting}[language=Rust]
// Anti-pattern: Primitive Obsession & Shotgun Validation
pub fn submit_international_wire(
    customer_id: u64,
    sender_account_id: u64,
    recipient_account_id: u64,
    destination_iban: &str,
    currency_iso: &str,
    amount_cents: i64,
    fee_cents: i64,
) -> Result<u64, String> {
    if sender_account_id == recipient_account_id { /* ... */ }
    if amount_cents <= 0 || fee_cents < 0 { /* ... */ }
    if currency_iso.len() != 3 { /* ... */ }
    if !is_valid_iban_checksum(destination_iban) { /* ... */ }
    // ...
    Ok(1)
}
\end{lstlisting}

Count the structural hazards hiding inside those seven primitive parameters:
\begin{enumerate}
  \item \textbf{Silent argument-swap bugs.} Three adjacent parameters have type \texttt{u64}, two have type \texttt{\&str}, and two have type \texttt{i64}. If a caller invokes \texttt{submit\_international\_wire(\allowbreak{}acct\_id,\allowbreak\ cust\_id,\allowbreak\ recip\_id,\allowbreak\ currency,\allowbreak\ iban,\allowbreak\ fee,\allowbreak\ amount)}, \texttt{rustc} sees \texttt{(u64,\allowbreak\ u64,\allowbreak\ u64,\allowbreak\ \&str,\allowbreak\ \&str,\allowbreak\ i64,\allowbreak\ i64)} --- a type-perfect match. In production, you have just debited the wrong account or swapped the transfer principal with the wire fee.
  \item \textbf{Shotgun validation.} Suppose \texttt{submit\_\allowbreak{}international\_\allowbreak{}wire} calls \texttt{is\_\allowbreak{}valid\_\allowbreak{}iban\_\allowbreak{}checksum} and passes \texttt{destination\_\allowbreak{}iban: \&str} down to \texttt{route\_\allowbreak{}correspondent\_\allowbreak{}bank} and \texttt{format\_\allowbreak{}swift\_\allowbreak{}mt103}. When you read \texttt{format\_\allowbreak{}swift\_\allowbreak{}mt103} six months later, how do you know its caller remembered to validate the IBAN? You don't. The type \texttt{\&str} includes \texttt{""}, \texttt{"hello\textbackslash nworld"}, and a 10-megabyte novel. So defensive programmers paste the exact same validation \texttt{if} checks at the top of every downstream helper --- scattering domain rules across dozens of files while paying repeated runtime CPU checks.
\end{enumerate}

\subsection{Validation Throws Proof Away; Parsing Preserves It in a Type}

In a classic 2019 essay originating in the Haskell community and adopted across the Rust ecosystem, Alexis King named the cure in three words: \textbf{Parse, don't validate}.

\begin{definition}[Validation versus Parsing]
Compare two function signatures operating on raw untrusted input:
\begin{align*}
  \text{Validate:} &\quad \texttt{\small fn validate\_iban(raw: \&str) -> Result<(), IbanError>} \\
  \text{Parse:}    &\quad \texttt{\small fn parse\_iban(raw: \&str) -> Result<Iban, IbanError>}
\end{align*}
Both inspect the exact same characters and reject the exact same malformed strings. Yet \texttt{validate\_\allowbreak{}iban} returns \texttt{()} --- the zero-byte unit type from Chapter~7 --- \emph{evaporating} the knowledge it just worked to establish and leaving the caller holding the same unrefined \texttt{\&str}. By contrast, \texttt{parse\_\allowbreak{}iban} narrows the output into a distinct type \texttt{Iban} whose only constructors enforce the invariant. Downstream functions accepting \texttt{iban:\allowbreak\ \&Iban} neither re-check the string nor return \texttt{Result} for malformed syntax; possessing an \texttt{Iban} value \emph{is} a machine-checked proof that parsing succeeded.
\end{definition}

Why didn't Java~6 developers wrap every \texttt{long} ID and positive cent amount in its own wrapper class (\texttt{new AccountId(42L)}, \texttt{new PositiveCents(5000L)})? Because on the JVM (prior to Project Valhalla value classes), every wrapper object carries a 12-to-16-byte object header, lives as a separate heap allocation behind an 8-byte pointer, can be \texttt{null}, stresses garbage-collection allocation rates, and turns arrays into arrays of pointers.

In Rust, wrapping a single field inside a tuple struct --- called the \textbf{Newtype pattern} --- has \textbf{zero runtime memory or CPU cost}.

\subsection{Zero-Cost Newtypes, \texttt{\#[repr(transparent)]}, and Module Privacy}

Let us build a real domain module \texttt{mod domain} containing distinct ID types, a strictly positive cent type \texttt{PositiveCents}, and an ISO~13616 \texttt{Iban} newtype that verifies country length, alphanumeric characters, and the international modulo-97 check-digit equation at construction time:

\begin{lstlisting}[language=Rust]
pub mod domain {
    use std::fmt;
    use std::num::NonZeroI64;
    use std::str::FromStr;

    // 1. Distinct zero-cost identifier newtypes (inner field is PRIVATE!)
    #[repr(transparent)]
    #[derive(Copy, Clone, Debug, PartialEq, Eq, Hash, PartialOrd, Ord)]
    pub struct CustomerId(u64);

    #[repr(transparent)]
    #[derive(Copy, Clone, Debug, PartialEq, Eq, Hash, PartialOrd, Ord)]
    pub struct AccountId(u64);

    impl CustomerId {
        pub const fn new(raw: u64) -> Self { Self(raw) }
        pub const fn get(self) -> u64 { self.0 }
    }

    impl AccountId {
        pub const fn new(raw: u64) -> Self { Self(raw) }
        pub const fn get(self) -> u64 { self.0 }
    }

    // 2. Strictly positive minor-unit money (> 0 cents guaranteed!)
    #[repr(transparent)]
    #[derive(Copy, Clone, Debug, PartialEq, Eq, PartialOrd, Ord)]
    pub struct PositiveCents(NonZeroI64);

    #[derive(Debug, Clone, PartialEq, Eq)]
    pub enum MoneyError {
        NonPositiveAmount(i64),
        Overflow,
    }

    impl PositiveCents {
        pub fn try_new(cents: i64) -> Result<Self, MoneyError> {
            if cents <= 0 {
                return Err(MoneyError::NonPositiveAmount(cents));
            }
            // Safe: we just proved cents > 0, so NonZeroI64::new is Some
            let nz = NonZeroI64::new(cents)
                .ok_or(MoneyError::NonPositiveAmount(cents))?;
            Ok(Self(nz))
        }

        pub const fn get(self) -> i64 {
            self.0.get()
        }

        #[must_use = "checked_add returns a new Result without mutating"]
        pub fn checked_add(self, rhs: Self) -> Result<Self, MoneyError> {
            let sum = self.get()
                .checked_add(rhs.get())
                .ok_or(MoneyError::Overflow)?;
            Self::try_new(sum)
        }
    }

    // 3. Validated IBAN newtype with ISO 13616 modulo-97 checksum
    #[derive(Clone, Debug, PartialEq, Eq, Hash)]
    pub struct Iban(String); // Private inner String!

    #[derive(Debug, Clone, PartialEq, Eq)]
    pub enum IbanError {
        InvalidLength(usize),
        InvalidCountryHeader,
        InvalidCharacter(char),
        ChecksumFailed,
    }

    impl FromStr for Iban {
        type Err = IbanError;

        fn from_str(raw: &str) -> Result<Self, Self::Err> {
            let canonical: String = raw
                .chars()
                .filter(|c| !c.is_ascii_whitespace())
                .map(|c| c.to_ascii_uppercase())
                .collect();

            if !(15..=34).contains(&canonical.len()) {
                return Err(IbanError::InvalidLength(canonical.len()));
            }
            let b = canonical.as_bytes();
            if !b[0].is_ascii_uppercase() || !b[1].is_ascii_uppercase()
                || !b[2].is_ascii_digit() || !b[3].is_ascii_digit()
            {
                return Err(IbanError::InvalidCountryHeader);
            }

            // ISO 13616: move first 4 chars to end, map A..Z -> 10..35,
            // and check that the integer has remainder 1 modulo 97.
            let rearranged =
                canonical[4..].chars().chain(canonical[..4].chars());
            let mut rem: u32 = 0;
            for ch in rearranged {
                if ch.is_ascii_digit() {
                    let digit = ch as u32 - b'0' as u32;
                    rem = (rem * 10 + digit).rem_euclid(97);
                } else if ch.is_ascii_uppercase() {
                    let val = (ch as u32 - b'A' as u32) + 10;
                    rem = (rem * 100 + val).rem_euclid(97);
                } else {
                    return Err(IbanError::InvalidCharacter(ch));
                }
            }
            if rem != 1 {
                return Err(IbanError::ChecksumFailed);
            }
            Ok(Self(canonical))
        }
    }

    impl Iban {
        pub fn as_str(&self) -> &str { &self.0 }
        pub fn country_code(&self) -> &str { &self.0[0..2] }
    }

    impl fmt::Display for Iban {
        fn fmt(&self, f: &mut fmt::Formatter<'_>) -> fmt::Result {
            f.write_str(&self.0)
        }
    }
}
\end{lstlisting}

Examine three architectural details working together in \texttt{mod domain}:

\begin{enumerate}
  \item \textbf{Field privacy is what gives a Newtype teeth.} Notice we declared \texttt{PositiveCents(\allowbreak{}NonZeroI64)} with a private inner field rather than \texttt{pub NonZeroI64}. Because the tuple field \texttt{.0} is private to \texttt{mod domain}, code outside \texttt{mod domain} cannot construct \texttt{PositiveCents(-500)} using literal syntax and cannot mutate the inner value directly. Every value of type \texttt{PositiveCents} in the entire program \emph{must} have passed through \texttt{PositiveCents::\allowbreak{}try\_new} or \texttt{checked\_add}. This pattern --- private fields paired with fallible factory functions returning \texttt{Result<Self, Err>} --- is called a \textbf{smart constructor}.
  \item \textbf{\texttt{\#[repr(transparent)]} and register-level zero cost.} Putting \texttt{\#[repr(transparent)]} on a single-field struct instructs the compiler that the wrapper must have the \emph{exact same size, alignment, and C/Rust calling-convention ABI} as its inner field. Look at what happens down at the hardware level:
\end{enumerate}

\begin{happeningbox}
On 64-bit x86-64 Linux/macOS (System V AMD64 ABI):
\begin{itemize}
  \item \texttt{size\_of::<AccountId>() == 8} and \texttt{align\_of::<AccountId>() == 8}. Passing \texttt{(CustomerId, AccountId, PositiveCents)} passes three plain 64-bit integers directly in CPU registers \texttt{rdi}, \texttt{rsi}, and \texttt{rdx} --- zero pointers, zero heap headers, zero indirection.
  \item Because \texttt{PositiveCents} wraps \texttt{std::num::NonZeroI64} (which forbids the bit pattern \texttt{0x0000\_0000\_0000\_0000}), the compiler uses \texttt{0} as the \emph{niche discriminant} for \texttt{None}: \texttt{size\_of::<Option<PositiveCents>>() == 8}. An optional validated positive dollar amount takes the exact same 8 bytes on the stack as a bare C \texttt{int64\_t}.
\end{itemize}
Once \texttt{rustc} finishes type-checking, \texttt{AccountId} and \texttt{CustomerId} erase completely from the emitted machine instructions. All the safety lives at compile time; zero bytes remain at runtime.
\end{happeningbox}

\begin{enumerate}\setcounter{enumi}{2}
  \item \textbf{Solving Chapter~6's Orphan Rule cleanly.} Recall from Chapter~6's coherence law that Rust forbids implementing a foreign trait (say \texttt{std::fmt::Display} or \texttt{sqlx::Encode}) for a foreign type (say \texttt{Vec<LedgerEntry>} or a third-party crate's \texttt{Decimal}). A Newtype \texttt{pub struct StatementReport(pub Vec<LedgerEntry>);} is a \emph{local} type owned by your crate, so implementing \texttt{fmt::Display for StatementReport} is 100-percent legal and keeps your formatting rules cleanly scoped without polluting global method resolution on every \texttt{Vec}.
\end{enumerate}

Look at our wire-submission signature now:

\begin{lstlisting}[language=Rust]
pub fn submit_international_wire(
    customer: domain::CustomerId,
    sender: domain::AccountId,
    recipient: domain::AccountId,
    destination: &domain::Iban,
    amount: domain::PositiveCents,
) -> Result<WireReceipt, WireError>
\end{lstlisting}

Swapping \texttt{customer} with \texttt{sender} is an immediate compiler error. Passing a negative or zero amount is impossible. Re-checking the IBAN checksum inside \texttt{submit\_international\_wire} is unnecessary because \texttt{\&Iban} already proves ISO~13616 validity. Validation happened once, at the HTTP JSON deserialization edge, and the type system carried the proof everywhere else for free.

% ======================================================================
\section{Compile-Time State Machines: The Typestate Pattern and Builders}
\label{ch08:typestate-and-builders}
% ======================================================================

Newtypes enforce invariants on \emph{individual values}. What about invariants across \emph{time} --- multi-step protocols where operations are only legal in a strict order?

Think about the lifecycle of a high-value wire order inside a bank:
\begin{enumerate}
  \item An API handler creates a \textbf{Draft} wire order holding the sender ID, recipient IBAN, and \texttt{PositiveCents} amount.
  \item A sanctions and anti-money-laundering (AML/OFAC) engine screens the counterparties. Only if screening passes does the order become \textbf{ComplianceCleared}, acquiring a cryptographic officer/screening stamp.
  \item A settlement worker transmits a cleared order to the central clearing rail, producing a \textbf{Settled} wire with a rail transaction sequence number.
\end{enumerate}

Two bugs must never occur: (a) settling a \texttt{Draft} wire that bypassed AML compliance screening, and (b) settling the \emph{same} \texttt{ComplianceCleared} wire twice.

In Java~6, a single mutable \texttt{WireOrder} class tries to police all three phases at runtime using nullable fields and boolean flags:

\begin{lstlisting}[language=Java]
// Java 6: temporal protocol encoded via mutable flags + runtime exceptions
public class WireOrder {
    private boolean complianceCleared = false;
    private boolean settled = false;
    private String ofacTicket; // null until complianceCleared == true!
    private Long railTxId;     // null until settled == true!

    public void approveCompliance(String ticket) {
        if (this.complianceCleared || this.settled) {
            throw new IllegalStateException("Invalid state for approval");
        }
        this.complianceCleared = true;
        this.ofacTicket = ticket;
    }

    public void settle(long txId) {
        if (!this.complianceCleared || this.settled) {
            throw new IllegalStateException(
                "Wire not cleared or already settled!"
            );
        }
        this.settled = true;
        this.railTxId = txId;
    }
}
\end{lstlisting}

Every caller can invoke \texttt{order.settle(999L)} on a fresh draft; the Java compiler says ``looks fine to me,'' and you only find out when an \texttt{IllegalStateException} is thrown in production (or worse, if someone forgets the \texttt{if (!this.complianceCleared)} guard inside a new batch-settlement method, an unscreened wire hits the SWIFT network). And notice those two nullable fields \texttt{ofacTicket} and \texttt{railTxId}: two booleans plus two nullable fields create $2 \times 2 \times 2 \times 2 = 16$ representable combinations in memory when only $3$ states are legal.

\subsection{Encoding Protocol States in Type Parameters and \texttt{PhantomData}}

In Rust, we can lift the state of the protocol out of runtime booleans and put it directly into a compile-time generic parameter \texttt{S} on \texttt{WireOrder<S>}. This is the \textbf{Typestate pattern}.

You actually have two complementary ways to represent the state parameter \texttt{S}, and you should know both:

\begin{enumerate}
  \item \textbf{Pure Zero-Sized Marker Types + \texttt{PhantomData<S>}.} When two states store the exact same runtime fields and differ only in which methods are permitted (for instance, \texttt{RawQuery} versus \texttt{SanitizedQuery}, or \texttt{Locked} versus \texttt{Unlocked}), we define empty unit structs \texttt{pub struct Draft;} and \texttt{pub struct Verified;} that occupy \textbf{0 bytes} in RAM. Because Rust requires every generic type parameter \texttt{S} declared on \texttt{struct Thing<S>} to be referenced in at least one field (so the compiler knows \texttt{S}'s variance and drop behavior from Chapter~4), we include a zero-byte marker field \texttt{\_state:\allowbreak\ std::marker::PhantomData<S>}.
  \item \textbf{State-Specific Payload Structs (\texttt{state: S}).} Even better, when entering a state \emph{produces new data} that did not exist before --- such as \texttt{ComplianceCleared} acquiring an \texttt{ofac\_ticket} and \texttt{Settled} acquiring a \texttt{rail\_tx\_id} --- we store \texttt{state: S} directly as a field. For \texttt{Draft}, \texttt{pub struct Draft;} is still a 0-byte zero-sized type (ZST), so \texttt{WireOrder<Draft>} pays zero bytes for \texttt{state}. For \texttt{ComplianceCleared}, \texttt{state} holds the \texttt{ApprovalStamp} directly without any \texttt{Option<...>} wrapper.
\end{enumerate}

\begin{formalbox}
\textbf{What \texttt{PhantomData<S>} Is at the Byte and Type Level.}
Inside \texttt{core::marker}, \texttt{PhantomData<T>} is a compiler-intrinsic zero-sized type:
\begin{align*}
  \texttt{std::mem::size\_of::<PhantomData<S>>()} &= 0, \\
  \texttt{std::mem::align\_of::<PhantomData<S>>()} &= 1.
\end{align*}
At runtime, a \texttt{PhantomData<S>} field emits \emph{zero bits} in memory and \emph{zero instructions} during moves or copies. At compile time, it tells \texttt{rustc} to treat the enclosing struct ``as if'' it logically owned a value of type \texttt{S} for type discrimination, lifetime variance, and \texttt{Send}/\texttt{Sync} auto-trait derivation. (If you want invariance without drop/send ownership effects, a common systems idiom is \texttt{PhantomData<fn(S) -> S>}, which is always \texttt{Copy + Send + Sync} regardless of \texttt{S}.)
\end{formalbox}

\subsection{Why Rust's Destructive Move Makes Typestate Airtight}

Here is the deeper insight that surprises engineers who have tried to emulate typestate using generics in Java, C++, or TypeScript. Suppose in Java you wrote \texttt{ClearedWire approve(DraftWire d)}. Nothing in Java stops the caller from writing:

\begin{lstlisting}[language=Java]
// In Java / TypeScript / Go, the old reference remains usable!
ClearedWire c1 = draft.approve("OFAC-OK-1");
ClearedWire c2 = draft.approve("OFAC-OK-2"); // Compiles! Double-approved!
\end{lstlisting}

Even in C++11, \texttt{std::move(draft)} leaves \texttt{draft} as a live lvalue on the stack in a ``valid but unspecified'' state whose methods can still be called by mistake.

In Rust, look at what happens when transition methods take \textbf{\texttt{self} by value} (not \texttt{\&self}, and not \texttt{\&mut self}):

\begin{lstlisting}[language=Rust]
use std::marker::PhantomData;
use crate::domain::{AccountId, Iban, PositiveCents};

// State types: Draft is 0 bytes; later states carry their proof artifacts!
pub struct Draft;

pub struct ComplianceCleared {
    pub officer_id: u32,
    pub ofac_ticket: [u8; 16],
}

pub struct Settled {
    pub officer_id: u32,
    pub ofac_ticket: [u8; 16],
    pub rail_tx_id: u64,
}

pub struct Rejected {
    pub reason: String,
}

// Generic wire order parameterized by its current protocol state S,
// plus a zero-byte PhantomData marker for the rail protocol R.
pub struct SwiftRail;

pub struct WireOrder<S, R = SwiftRail> {
    sender: AccountId,
    destination: Iban,
    amount: PositiveCents,
    state: S,
    _rail: PhantomData<R>, // 0 bytes at runtime!
}

// Phase 1 methods: ONLY exist on WireOrder<Draft, R>!
impl<R> WireOrder<Draft, R> {
    pub fn new_draft(
        sender: AccountId,
        destination: Iban,
        amount: PositiveCents,
    ) -> Self {
        Self {
            sender,
            destination,
            amount,
            state: Draft,
            _rail: PhantomData,
        }
    }

    // Takes `self` BY VALUE: consumes WireOrder<Draft, R> permanently!
    pub fn approve_compliance(
        self,
        officer_id: u32,
        ofac_ticket: [u8; 16],
    ) -> WireOrder<ComplianceCleared, R> {
        WireOrder {
            sender: self.sender,
            destination: self.destination,
            amount: self.amount,
            state: ComplianceCleared { officer_id, ofac_ticket },
            _rail: PhantomData,
        }
    }

    pub fn reject(
        self,
        reason: impl Into<String>,
    ) -> WireOrder<Rejected, R> {
        WireOrder {
            sender: self.sender,
            destination: self.destination,
            amount: self.amount,
            state: Rejected { reason: reason.into() },
            _rail: PhantomData,
        }
    }
}

// Phase 2 methods: ONLY exist on WireOrder<ComplianceCleared, R>!
impl<R> WireOrder<ComplianceCleared, R> {
    pub fn ofac_ticket(&self) -> &[u8; 16] {
        &self.state.ofac_ticket
    }

    // Consumes WireOrder<ComplianceCleared, R> -> cannot settle twice!
    pub fn settle(self, rail_tx_id: u64) -> WireOrder<Settled, R> {
        WireOrder {
            sender: self.sender,
            destination: self.destination,
            amount: self.amount,
            state: Settled {
                officer_id: self.state.officer_id,
                ofac_ticket: self.state.ofac_ticket,
                rail_tx_id,
            },
            _rail: PhantomData,
        }
    }
}

// Shared read-only accessors available in EVERY state S:
impl<S, R> WireOrder<S, R> {
    pub fn sender(&self) -> AccountId { self.sender }
    pub fn destination(&self) -> &Iban { &self.destination }
    pub fn amount(&self) -> PositiveCents { self.amount }
}
\end{lstlisting}

Try breaking this protocol in client code:
\begin{itemize}
  \item \textbf{Try settling an unscreened draft:} \texttt{draft.settle(1001)} fails to compile with \texttt{no method named `settle` found for struct `WireOrder<Draft>`}. The method does not exist in \texttt{Draft}'s method table.
  \item \textbf{Try settling a cleared wire twice:}
\begin{lstlisting}[language=Rust]
let cleared = draft.approve_compliance(7, *b"OFAC-2026-PASS01");
let done_1 = cleared.settle(90001); // `cleared` is MOVED here!
// let done_2 = cleared.settle(90002);
// ^-- COMPILE ERROR: borrow/use of moved value: `cleared`!
\end{lstlisting}
\end{itemize}

Because \texttt{WireOrder} deliberately does \emph{not} derive \texttt{Clone}, and \texttt{settle(self, ...)} takes ownership by value, Rust's affine ownership system enforces \emph{linear state progression}: one \texttt{Draft} token enters, one \texttt{ComplianceCleared} token emerges, and the \texttt{Draft} binding ceases to exist in the compiler's live-variable table.

\begin{figure}[H]
\centering
\begin{tikzpicture}[
    font=\small,
    state/.style={draw=sourceblue!85!black, fill=sourceblue!6, rounded corners=4pt, very thick, minimum width=3.45cm, minimum height=1.35cm, align=center},
    termstate/.style={draw=targetgreen!80!black, fill=targetgreen!7, rounded corners=4pt, very thick, minimum width=3.25cm, minimum height=1.25cm, align=center},
    rejstate/.style={draw=quotientorange!85!black, fill=quotientorange!8, rounded corners=4pt, very thick, minimum width=3.25cm, minimum height=1.15cm, align=center},
    trans/.style={-{Stealth[length=2.6mm,width=1.9mm]}, very thick, draw=sourceblue!85!black},
    illegal/.style={-{Stealth[length=2.4mm,width=1.8mm]}, thick, dashed, draw=trapred!85!black}
]
  % States
  \node[state] (draft) at (1.85,2.75) {
    \textbf{\texttt{WireOrder<Draft>}}\\[2pt]
    \scriptsize\texttt{state: Draft} (0 B)\\[-1pt]
    \scriptsize\texttt{\_rail: PhantomData} (0 B)
  };

  \node[state] (cleared) at (6.95,2.75) {
    \textbf{\texttt{WireOrder<Cleared>}}\\[2pt]
    \scriptsize\texttt{officer\_id: u32}\\[-1pt]
    \scriptsize\texttt{ofac\_ticket: [u8; 16]}
  };

  \node[termstate] (settled) at (11.95,3.55) {
    \textbf{\texttt{WireOrder<Settled>}}\\[2pt]
    \scriptsize\texttt{+ rail\_tx\_id: u64}
  };

  \node[rejstate] (rejected) at (11.95,1.55) {
    \textbf{\texttt{WireOrder<Rejected>}}\\[2pt]
    \scriptsize\texttt{reason: String}
  };

  % Legal consuming transitions
  \draw[trans] (draft.east) -- node[above, font=\scriptsize\ttfamily] {approve(self)} node[below, font=\tiny, text=sourceblue!85!black] {moves \texttt{self}} (cleared.west);
  \draw[trans, draw=targetgreen!80!black] (cleared.east) -- ++(0.35,0) |- node[pos=0.72, above, font=\scriptsize\ttfamily, text=targetgreen!75!black] {settle(self)} (settled.west);
  \draw[trans, draw=quotientorange!85!black] (draft.south) |- node[pos=0.68, below, font=\scriptsize\ttfamily, text=quotientorange!85!black] {reject(self) (consumes Draft)} (rejected.west);

  % Illegal transitions caught at compile time
  \draw[illegal, bend left=27] (draft.north east) to node[above, font=\scriptsize, text=trapred!90!black] {\textbf{$\times$} \texttt{draft.settle()} --- no such method on \texttt{Draft}} (settled.north west);

  \draw[illegal] (cleared.north) .. controls +( -0.8, 0.95) and +( 0.8, 0.95) .. node[above, font=\scriptsize, text=trapred!90!black] {\textbf{$\times$} 2nd \texttt{settle()}: use of moved value} (cleared.north east);
\end{tikzpicture}
\caption{Compile-time Typestate automaton for \texttt{WireOrder<S>}. Valid transitions consume \texttt{self} by value; out-of-order calls (\texttt{draft.settle()}) and duplicate settlements fail at compile time with zero runtime flag overhead.}
\label{ch08:fig-typestate-automaton}
\end{figure}

\subsection{Compile-Time Typestate vs Runtime \texttt{enum} State Machines}

Once engineers discover Typestate, a common temptation is to try to use it for \emph{every} state machine in a codebase --- and promptly hit a wall the first time they load rows from PostgreSQL.

Ask yourself: if you run \texttt{SELECT id, status, sender\_id, iban, cents FROM wire\_orders WHERE customer\_id = \$1}, what is the return type of that function? At compile time, \texttt{rustc} cannot possibly know whether row~0 in PostgreSQL currently has \texttt{status = 'DRAFT'} or \texttt{status = 'SETTLED'}. Even more directly, a \texttt{Vec<T>} requires every element to have the \emph{exact same concrete type} \texttt{T}, so you cannot put a \texttt{WireOrder<Draft>} and a \texttt{WireOrder<Settled>} side by side in the same \texttt{Vec}.

Here is the crisp architectural decision rule separating Chapter~5's runtime \texttt{enum}s from Chapter~8's compile-time Typestate --- and how production Rust combines both:

\begin{intuitionbox}
\begin{itemize}
  \item \textbf{Use Compile-Time Typestate (\texttt{Order<S>})} when the state progression is \emph{statically choreographed by the code path itself}: builders, multi-step cryptographic handshakes, SQL transaction lifecycles, or a handler pipeline that always parses a request into \texttt{Draft}, screens it into \texttt{ComplianceCleared}, and settles it.
  \item \textbf{Use a Runtime Sum Type (\texttt{enum AnyWireOrder})} when state depends on \emph{external data unknown until runtime} (database rows, network messages, user events) or when you must store a \emph{heterogeneous collection} (\texttt{Vec<AnyWireOrder>}) of items currently sitting at different lifecycle stages.
\end{itemize}
\end{intuitionbox}

You can combine the two. Wrap your strongly typed typestate structs \emph{inside} a runtime \texttt{enum} envelope at storage and network boundaries:

\begin{lstlisting}[language=Rust]
// Runtime envelope wrapping compile-time typestate structs!
pub enum AnyWireOrder {
    Draft(WireOrder<Draft>),
    Cleared(WireOrder<ComplianceCleared>),
    Settled(WireOrder<Settled>),
    Rejected(WireOrder<Rejected>),
}

pub fn process_pending_from_db(
    row: AnyWireOrder,
    officer_id: u32,
) -> AnyWireOrder {
    match row {
        // Unpacking `Draft(draft)` yields a typed WireOrder<Draft>!
        AnyWireOrder::Draft(draft) => {
            let cleared =
                draft.approve_compliance(officer_id, *b"OFAC-AUTO-OK-001");
            AnyWireOrder::Cleared(cleared)
        }
        // Cleared, settled, or rejected rows cannot be re-approved:
        other => other,
    }
}
\end{lstlisting}

Look how cleanly the two layers interlock: PostgreSQL hydration produces \texttt{AnyWireOrder}, \texttt{match} narrows a dynamic variant into a concrete \texttt{WireOrder<Draft>}, and inside that branch all of Typestate's compile-time guarantees govern every transition.

\subsection{Idiomatic Builders: Owned Chaining and Typestate Builders}

Unlike Java~6, Rust intentionally has neither method overloading (you cannot write four \texttt{new(...)} constructors with 2, 3, 5, and 8 parameters) nor default/optional function arguments. Whenever a configuration or request struct has more than three or four knobs, idiomatic Rust provides a \textbf{Builder}.

Two flavors appear across the ecosystem:

\paragraph{1. Owned-Chaining Builder (for optional knobs with sensible defaults).}
When all required arguments fit cleanly in \texttt{new(...)} and the remaining settings have safe defaults, each builder method takes \texttt{mut self} by value, updates one field, and returns \texttt{Self}:

\begin{lstlisting}[language=Rust]
pub struct RetryPolicy {
    max_attempts: u32,
    initial_backoff_ms: u64,
    max_backoff_ms: u64,
    jitter: bool,
}

impl RetryPolicy {
    #[must_use]
    pub fn new() -> Self {
        Self {
            max_attempts: 3,
            initial_backoff_ms: 50,
            max_backoff_ms: 2_000,
            jitter: true,
        }
    }

    #[must_use]
    pub fn max_attempts(mut self, attempts: u32) -> Self {
        self.max_attempts = attempts.max(1);
        self
    }

    #[must_use]
    pub fn backoff_bounds_ms(mut self, initial: u64, max: u64) -> Self {
        self.initial_backoff_ms = initial;
        self.max_backoff_ms = max.max(initial);
        self
    }

    #[must_use]
    pub fn jitter(mut self, enabled: bool) -> Self {
        self.jitter = enabled;
        self
    }
}
\end{lstlisting}

Taking \texttt{mut self -> Self} supports both fluent chains like \texttt{Retry\allowbreak{}Policy::\allowbreak{}new().\allowbreak{}max\_\allowbreak{}attempts(5).\allowbreak{}jitter(false)} \emph{and} conditional updates inside loops (\texttt{policy = policy.\allowbreak{}max\_\allowbreak{}attempts(10)}) without dangling temporary \texttt{\&mut} borrows.

\paragraph{2. Typestate Builder (for mandatory configuration with zero \texttt{.unwrap()}).}
What if a \texttt{LedgerPoolBuilder} has \emph{mandatory} fields without sensible defaults --- say \texttt{postgres\_url} and \texttt{signing\_key} --- alongside optional timeouts? A Java~6 builder initializes those fields to \texttt{null} and throws \texttt{IllegalStateException("missing postgresUrl")} inside \texttt{.build()}. Using Typestate, \texttt{.build()} is only implemented on \texttt{LedgerPoolBuilder<HasUrl, HasKey>}. Forgetting to call \texttt{.postgres\_url(...)} makes \texttt{.build()} a compile-time error rather than a startup crash:

\begin{lstlisting}[language=Rust]
pub struct Missing;
pub struct Provided<T>(T);

pub struct LedgerPoolBuilder<U, K> {
    url: U,
    signing_key: K,
    max_connections: u32,
}

impl LedgerPoolBuilder<Missing, Missing> {
    pub fn new() -> Self {
        Self { url: Missing, signing_key: Missing, max_connections: 16 }
    }
}

impl<U, K> LedgerPoolBuilder<U, K> {
    pub fn max_connections(mut self, n: u32) -> Self {
        self.max_connections = n.max(1);
        self
    }

    pub fn postgres_url(
        self,
        url: impl Into<String>,
    ) -> LedgerPoolBuilder<Provided<String>, K> {
        LedgerPoolBuilder {
            url: Provided(url.into()),
            signing_key: self.signing_key,
            max_connections: self.max_connections,
        }
    }

    pub fn signing_key(
        self,
        key: [u8; 32],
    ) -> LedgerPoolBuilder<U, Provided<[u8; 32]>> {
        LedgerPoolBuilder {
            url: self.url,
            signing_key: Provided(key),
            max_connections: self.max_connections,
        }
    }
}

pub struct LedgerPoolConfig {
    pub url: String,
    pub signing_key: [u8; 32],
    pub max_connections: u32,
}

// .build() ONLY exists once BOTH mandatory fields have been supplied!
impl LedgerPoolBuilder<Provided<String>, Provided<[u8; 32]>> {
    pub fn build(self) -> LedgerPoolConfig {
        LedgerPoolConfig {
            url: self.url.0,
            signing_key: self.signing_key.0,
            max_connections: self.max_connections,
        }
    }
}
\end{lstlisting}

Notice that the caller can supply \texttt{.signing\_key(...)} and \texttt{.postgres\_url(...)} in \emph{either} order, interleave \texttt{.max\_connections(32)} anywhere, and \texttt{.build()} returns an infallible \texttt{LedgerPoolConfig} directly --- no \texttt{Option}, no \texttt{Result}, no runtime \texttt{.unwrap()}.

% ======================================================================
\section{Ergonomic, Zero-Copy API Boundaries and Borrow-Splitting Idioms}
\label{ch08:api-boundaries}
% ======================================================================

Once your data structures and state machines are cleanly modeled, the next place ``Java-in-Rust syndrome'' leaks CPU cycles and frustrates callers is at \emph{function boundaries} --- specifically, how functions accept borrowed views, how constructors accept owned buffers, and how methods borrow fields of \texttt{self}.

\subsection{Why \texttt{\&String}, \texttt{\&Vec<T>}, and \texttt{\&Box<T>} Are Anti-Patterns}

A developer accustomed to Java's \texttt{String} and \texttt{List<Posting>} (or C++'s \texttt{const std::string\&} and \texttt{const std::vector<Posting>\&}) naturally writes Rust helper signatures like these:

\begin{lstlisting}[language=Rust]
// Anti-pattern: borrowing the heap-owning container header!
pub fn is_sepa_country(iban: &String) -> bool {
    matches!(&iban[0..2], "DE" | "FR" | "NL" | "ES" | "IT" | "CH")
}

pub fn net_balance_delta(postings: &Vec<i64>) -> i64 {
    postings.iter().copied().sum()
}

pub fn audit_fee_rule(rule: &Box<FeeSchedule>) -> i64 {
    rule.base_fee_cents
}
\end{lstlisting}

Every seasoned Rust reviewer flags \texttt{\&String}, \texttt{\&Vec<T>}, and \texttt{\&Box<T>} immediately. Look at what is wrong on \emph{both} the caller side and the hardware side:

\begin{enumerate}
  \item \textbf{You lock out valid callers who do not own a \texttt{String} or \texttt{Vec}.} Suppose a caller wants to check \texttt{is\_sepa\_country(\allowbreak{}"DE89370400440532013000")} using a string literal in read-only \texttt{.rodata} memory (\texttt{\&'static str}), or slice an IBAN field directly out of a 4-kilobyte network socket buffer (\texttt{\&packet\_str[8..30]}). Neither caller possesses a heap-allocated \texttt{String}. Forcing \texttt{\&String} makes them call \texttt{.to\_string()} just to allocate a throwaway heap copy so your read-only function can glance at two bytes and immediately drop it. Similarly, \texttt{\&Vec<i64>} rejects a fixed-size stack array \texttt{[100, -50, -50]} or a window slice \texttt{\&all\_postings[10..20]}.
  \item \textbf{You introduce a pointless double pointer dereference.} Recall the stack layout from Chapter~3: a \texttt{String} on the stack is 24 bytes \texttt{(ptr, len, cap)}. Passing \texttt{iban: \&String} passes an 8-byte pointer to that 24-byte stack \texttt{String}, which the CPU must first dereference just to load \texttt{ptr} and \texttt{len} before dereferencing \texttt{ptr} a second time to read the UTF-8 characters on the heap.
\end{enumerate}

The idiomatic replacement is universal: \textbf{whenever a function only reads data without growing or shrinking its backing allocation, accept the borrowed unsized view}:
\begin{itemize}
  \item Replace \texttt{\&String} with \textbf{\texttt{\&str}} (a 16-byte \texttt{(ptr, len)} fat pointer passed in two CPU registers, pointing directly at UTF-8 bytes whether they live on the heap, on the stack, or in static \texttt{.rodata}).
  \item Replace \texttt{\&Vec<T>} with \textbf{\texttt{\&[T]}} (a 16-byte \texttt{(ptr, len)} slice view over any contiguous \texttt{T} sequence).
  \item Replace \texttt{\&Box<T>} (or \texttt{\&Rc<T>}, \texttt{\&Arc<T>}) with \textbf{\texttt{\&T}}.
  \item For filesystem paths, accept \textbf{\texttt{path:\allowbreak\ impl AsRef<Path>}} (from \texttt{std::path::Path}) so callers can pass \texttt{"config/fees.toml"} (\texttt{\&str}), \texttt{PathBuf}, or \texttt{\&Path} without conversion ceremony.
\end{itemize}

And what about callers who \emph{do} already hold a \texttt{s: String} or \texttt{v: Vec<i64>}? Thanks to Chapter~6's \textbf{\texttt{Deref} coercion} (\texttt{String: Deref<Target = str>} and \texttt{Vec<T>: Deref<Target = [T]>}), writing \texttt{is\_sepa\_country(\&s)} and \texttt{net\_balance\_delta(\&v)} still compiles with zero changes at the call site. Switching from \texttt{\&String} to \texttt{\&str} strictly widens who can call your function while eliminating a pointer hop.

\subsection{Caller-Controlled Allocation: \texttt{impl Into<String>} and \texttt{Cow<'a, str>}}

Now consider the opposite case: what if a function or constructor \emph{does} need to store an owned \texttt{String} inside a struct?

\begin{lstlisting}[language=Rust]
pub struct AuditNote {
    actor: String,
    memo: String,
}

// Anti-pattern: taking &str when you ALWAYS allocate an owned String!
impl AuditNote {
    pub fn new_wasteful(actor: &str, memo: &str) -> Self {
        Self {
            actor: actor.to_owned(), // ALWAYS allocates + copies!
            memo: memo.to_owned(),
        }
    }

    // Idiomatic: let the caller move an existing String in O(1),
    // or allocate from &str only when needed!
    pub fn new(actor: impl Into<String>, memo: impl Into<String>) -> Self {
        Self {
            actor: actor.into(),
            memo: memo.into(),
        }
    }
}
\end{lstlisting}

Compare what happens when an HTTP JSON parser has \emph{already} allocated \texttt{parsed\_memo: String} on the heap:
\begin{itemize}
  \item Calling \texttt{AuditNote::\allowbreak{}new\_wasteful(\allowbreak{}\&parsed\_memo)} borrows the heap string as \texttt{\&str}, calls \texttt{malloc} to allocate a \emph{second} identical heap buffer, \texttt{memcpy}s the bytes, and frees \texttt{parsed\_memo} at the end of the caller's scope.
  \item Calling \texttt{AuditNote::\allowbreak{}new(actor\_literal,\allowbreak\ parsed\_memo)} monomorphizes two conversions: for \texttt{actor\_literal:\allowbreak\ \&str}, \texttt{.into()} allocates once; for \texttt{parsed\_memo:\allowbreak\ String}, \texttt{Into<String>::\allowbreak{}into(self)} is the identity function --- it moves 24 bytes on the stack in $O(1)$ time with \textbf{zero heap allocations}. (For tiny non-generic compile speed, accepting \texttt{memo: String} directly is equally idiomatic and makes allocation visible right at the call site.)
\end{itemize}

What about the trickiest middle ground of all: a function that \textbf{99 percent of the time can return a slice of its input untouched, and 1 percent of the time needs to allocate a modified string}?

In payment processing, every incoming ACH or SWIFT remittance memo must be normalized before ledger ingestion: leading/trailing ASCII whitespace is trimmed, and if the memo contains ASCII control characters (like \texttt{\textbackslash r}, \texttt{\textbackslash n}, or \texttt{\textbackslash t}), each control character must be replaced by a single space \texttt{' '}. In production traffic, 99 percent of bank memos are already clean printable ASCII. Returning \texttt{String} would force \texttt{malloc + memcpy} on all 100 percent of requests just to handle the dirty 1 percent.

Rust's standard library solves this with \textbf{\texttt{std::borrow::Cow<'a, B>}} (\emph{Clone-on-Write}):

\begin{lstlisting}[language=Rust]
// Simplified definition from std::borrow:
pub enum Cow<'a, B: ?Sized + ToOwned> {
    Borrowed(&'a B),               // Zero heap allocation: borrows input!
    Owned(<B as ToOwned>::Owned),  // Heap owned: e.g., String for B = str
}
\end{lstlisting}

Because \texttt{Cow<'a, str>} implements \texttt{Deref<Target = str>}, callers can read a \texttt{Cow<'a, str>} just like any \texttt{\&str} regardless of which variant came back, or call \texttt{.into\_owned()} if they later decide to persist an owned \texttt{String} (which moves the existing \texttt{String} out of \texttt{Cow::Owned} without re-allocating, or allocates once if it was \texttt{Cow::Borrowed}).

Here is \texttt{sanitize\_remittance\_memo} handling both paths:

\begin{lstlisting}[language=Rust]
use std::borrow::Cow;

/// Trims ASCII whitespace and replaces ASCII control chars with ' '.
/// Fast path (clean input): 0 heap allocations!
pub fn sanitize_remittance_memo(raw: &str) -> Cow<'_, str> {
    let trimmed = raw.trim_matches(|c: char| c.is_ascii_whitespace());

    // Scan once for the first offending control character
    let Some(first_bad_byte) = trimmed
        .bytes()
        .position(|b| b.is_ascii_control())
    else {
        // 99-percent fast path: return a 16-byte slice into `raw`!
        return Cow::Borrowed(trimmed);
    };

    // 1-percent slow path: allocate ONCE, copy clean prefix, sanitize rest
    let mut cleaned = String::with_capacity(trimmed.len());
    cleaned.push_str(&trimmed[..first_bad_byte]);
    for ch in trimmed[first_bad_byte..].chars() {
        if ch.is_ascii_control() {
            cleaned.push(' ');
        } else {
            cleaned.push(ch);
        }
    }
    Cow::Owned(cleaned)
}
\end{lstlisting}

\begin{figure}[H]
\centering
\begin{tikzpicture}[
    font=\small,
    mem/.style={draw, rounded corners=2pt, thick, align=left},
    arr/.style={-{Stealth[length=2.3mm,width=1.7mm]}, thick}
]
  % Left half: &str universal borrowed view
  \draw[draw=sourceblue!75!black, fill=sourceblue!4, rounded corners=4pt, thick] (0,0) rectangle (6.65,5.35);
  \node[font=\bfseries\small, text=sourceblue!85!black] at (3.32,4.98) {Universal View: \texttt{fn check(s: \&str)}};

  \node[mem, draw=structurepurple, fill=white, font=\scriptsize\ttfamily] (ro) at (1.85,4.05) {.rodata: "DE893704..."};
  \node[mem, draw=sourceblue!85!black, fill=white, font=\scriptsize\ttfamily] (hp) at (1.85,2.75) {Heap String buf};
  \node[mem, draw=targetgreen!80!black, fill=white, font=\scriptsize\ttfamily] (pk) at (1.85,1.45) {Socket buf[8..30]};

  \node[mem, draw=sourceblue!90!black, fill=white, very thick, font=\scriptsize\ttfamily, align=center] (slice) at (5.15,2.75) {
    \&str (16 B)\\
    ptr: *const u8\\
    len: usize
  };

  \draw[arr, structurepurple] (ro.east) -- (slice.north west);
  \draw[arr, sourceblue!85!black] (hp.east) -- (slice.west);
  \draw[arr, targetgreen!80!black] (pk.east) -- (slice.south west);
  \node[font=\scriptsize, text=sourceblue!85!black] at (3.32,0.42) {1 fat pointer accepts all 3 sources with 0 copies};

  % Right half: Cow<'a, str> zero-alloc fast path vs owned slow path
  \draw[draw=targetgreen!80!black, fill=targetgreen!4, rounded corners=4pt, thick] (7.05,0) rectangle (13.8,5.35);
  \node[font=\bfseries\small, text=targetgreen!75!black] at (10.42,4.98) {\texttt{Cow<'a, str>} (24 B Tagged Sum)};

  \node[mem, draw=sourceblue!80!black, fill=white, font=\scriptsize\ttfamily] (inbuf) at (10.42,4.15) {Caller input \&'a str: "INV-2026-0042"};

  \node[mem, draw=targetgreen!85!black, fill=white, thick, font=\scriptsize\ttfamily] (borrowed) at (8.85,2.35) {
    Cow::Borrowed\\
    ptr -> caller buf\\
    len: 13 (0 allocs!)
  };

  \node[mem, draw=quotientorange!90!black, fill=white, thick, font=\scriptsize\ttfamily] (owned) at (12.05,2.35) {
    Cow::Owned\\
    String \{ptr,len,cap\}\\
    (allocates on dirty)
  };

  \draw[arr, targetgreen!80!black] (borrowed.north) -- node[right, pos=0.45, font=\tiny, text=targetgreen!75!black] {clean input} (inbuf.south west);
  \node[mem, draw=quotientorange!90!black, fill=quotientorange!8, font=\scriptsize\ttfamily] (newheap) at (12.05,0.92) {New Heap Buffer};
  \draw[arr, quotientorange!90!black] (owned.south) -- (newheap.north);
  \node[font=\scriptsize, text=targetgreen!75!black] at (8.75,0.92) {Zero \texttt{malloc}!};
\end{tikzpicture}
\caption{Zero-copy boundary idioms. Left: accepting \texttt{\&str} (or \texttt{\&[T]}) lets static literals, heap \texttt{String}s, and network sub-slices enter through a single 16-byte register pair. Right: \texttt{Cow<'a, str>} avoids heap allocation on clean inputs while retaining the ability to allocate when normalization mutates bytes.}
\label{ch08:fig-slice-views-and-cow}
\end{figure}

\subsection{Defeating the Whole-Struct \texttt{\&mut self} Trap (Split Borrows)}

This is one of the most common borrow-checker errors in application code. Once you understand \emph{why} it happens, the fixes below are mechanical.

Suppose we build a \texttt{SettlementEngine} owning three distinct components: an in-memory \texttt{LedgerBook} of account balances, an immutable \texttt{FeeSchedule} loaded at startup, and an append-only \texttt{audit\_log}:

\begin{lstlisting}[language=Rust]
pub struct FeeSchedule {
    pub flat_wire_fee_cents: i64,
}

pub struct LedgerBook {
    pub balances: std::collections::HashMap<u64, i64>,
}

pub struct SettlementEngine {
    pub ledger: LedgerBook,
    pub fee_schedule: FeeSchedule,
    pub audit_log: Vec<String>,
}
\end{lstlisting}

Following OOP habits, you write a helper \texttt{apply\_fee(\&mut self, ...)} on \texttt{SettlementEngine} and invoke it inside \texttt{settle\_account(\&mut self)} while passing \texttt{\&self.\allowbreak{}fee\_schedule}:

\begin{lstlisting}[language=Rust]
impl SettlementEngine {
    // Helper mutates ONLY self.ledger and self.audit_log...
    // ...yet its signature demands `&mut self` (the ENTIRE struct)!
    fn apply_fee(&mut self, acct: u64, rules: &FeeSchedule) {
        let bal = self.ledger.balances.entry(acct).or_insert(0);
        *bal -= rules.flat_wire_fee_cents;
        self.audit_log.push(format!("charged fee to {acct}"));
    }

    pub fn settle_account(&mut self, acct: u64) {
        // COMPILE ERROR E0502: cannot borrow `self` as mutable because
        // `self.fee_schedule` is also borrowed as immutable!
        // self.apply_fee(acct, &self.fee_schedule);
    }
}
\end{lstlisting}

The body of \texttt{apply\_fee} only mutates the ledger and the audit log; it never touches the fee schedule. Why can't \texttt{rustc} just peek inside \texttt{apply\_fee} and see that the fields are disjoint?

\begin{whybox}
\textbf{Signature Modularity Is Deliberate.} In Rust, borrow checking is strictly \emph{modular across function boundaries}: \texttt{rustc} checks a caller against the callee's \textbf{signature alone}, never its body. Why? Suppose \texttt{apply\_fee(\&mut self, rules: \&FeeSchedule)} were allowed today because its body didn't mutate \texttt{self.fee\_schedule}. Six months later, a teammate adds \texttt{self.fee\_schedule.flat\_wire\_fee\_cents += 1;} inside \texttt{apply\_fee}. That edit would silently create an active mutable alias against \texttt{rules: \&FeeSchedule} \emph{or} break distant callers across the codebase without the method signature ever changing. Writing \texttt{\&mut self} is a contract claiming exclusive access to \textbf{every field} of \texttt{Self}.
\end{whybox}

Inside a \emph{single} function body, however, \texttt{rustc}'s borrow checker \emph{does} track every struct field independently. It knows \texttt{\&mut self.ledger}, \texttt{\&self.fee\_schedule}, and \texttt{\&mut self.audit\_log} occupy non-overlapping byte offsets inside \texttt{SettlementEngine}. That immediately gives us \textbf{three clean idioms} to cure whole-struct \texttt{\&mut self} conflicts:

\paragraph{Idiom 1: Put the method on the sub-struct it actually mutates (or make it an associated function taking disjoint fields).}
Instead of making \texttt{SettlementEngine} a ``God object'' where every helper takes \texttt{\&mut self}, put \texttt{apply\_fee} on \texttt{LedgerBook} (or pass \texttt{\&mut self.ledger} and \texttt{\&mut self.audit\_log} explicitly):

\begin{lstlisting}[language=Rust]
impl LedgerBook {
    pub fn apply_fee(
        &mut self,
        acct: u64,
        rules: &FeeSchedule,
        audit_log: &mut Vec<String>,
    ) {
        let bal = self.balances.entry(acct).or_insert(0);
        *bal -= rules.flat_wire_fee_cents;
        audit_log.push(format!("charged fee to {acct}"));
    }
}

impl SettlementEngine {
    pub fn settle_account(&mut self, acct: u64) {
        // Compiles cleanly! rustc sees three disjoint field borrows:
        //   &mut self.ledger, &self.fee_schedule, &mut self.audit_log
        self.ledger.apply_fee(
            acct,
            &self.fee_schedule,
            &mut self.audit_log,
        );
    }
}
\end{lstlisting}

\paragraph{Idiom 2: Group fields by mutation lifecycle.}
If five fields always change together while three configuration fields stay read-only after startup, group them into \texttt{state: EngineState} and \texttt{config: EngineConfig}. Then \texttt{self.state.step(\&self.config)} reads like plain English and never collides.

\paragraph{Idiom 3: Explicit ``View Structs'' (Borrow Projections) for complex subsystems.}
When several helpers need mutable access to two fields and shared access to two others, passing four separate parameters gets noisy. Define a lightweight \textbf{View Struct} holding only field references, plus a one-line \texttt{view\_mut(\&mut self)} projection method:

\begin{lstlisting}[language=Rust]
struct PostingView<'a> {
    ledger: &'a mut LedgerBook,
    audit_log: &'a mut Vec<String>,
    fees: &'a FeeSchedule,
}

impl<'a> PostingView<'a> {
    fn apply_fee(&mut self, acct: u64) {
        let bal = self.ledger.balances.entry(acct).or_insert(0);
        *bal -= self.fees.flat_wire_fee_cents;
        self.audit_log.push(format!("charged fee to {acct}"));
    }
}

impl SettlementEngine {
    fn posting_view(&mut self) -> PostingView<'_> {
        PostingView {
            ledger: &mut self.ledger,
            audit_log: &mut self.audit_log,
            fees: &self.fee_schedule,
        }
    }

    pub fn settle_batch(&mut self, accounts: &[u64]) {
        let mut view = self.posting_view();
        for &acct in accounts {
            view.apply_fee(acct);
        }
    }
}
\end{lstlisting}

Look at what \texttt{PostingView<'a>} achieves: the signature itself documents \emph{exact field-level permissions} (mutable ledger and audit log, read-only fee schedule), and every helper on \texttt{PostingView} enjoys clean \texttt{\&mut self} method syntax with zero borrow conflicts.

% ======================================================================
\section{Idiomatic Error Handling: \texttt{thiserror} vs \texttt{anyhow}, and \texttt{Result} vs \texttt{panic!}}
\label{ch08:error-design}
% ======================================================================

In Chapter~5 we saw how \texttt{Result<T, E>} and the \texttt{?} operator replace both C's silent integer error codes and Java~6's stack-unwinding exceptions. Yet knowing how \texttt{Result<T, E>} works mechanically leaves open an architectural question that trips up many new Rust teams: \emph{what should the error type \texttt{E} actually be across a multi-layer codebase?}

If you make \texttt{E = String} everywhere, callers cannot \texttt{match} on failure causes without fragile substring searches. Conversely, if you force an ops CLI migration script to define a 25-variant custom \texttt{enum} just to print an error message and exit with code~1, you drown in ceremony.

The Rust ecosystem converged on a two-zone rule based on a single question: \textbf{will the caller programmatically branch on which failure occurred, or will it report/log the failure up the stack?}

\subsection{Zone 1: Domain and Library Layers Use Structured \texttt{enum}s (\texttt{thiserror})}

Inside your domain logic, database repositories, and reusable library crates, callers \emph{act differently} depending on what failed:
\begin{itemize}
  \item In our Part~V Axum HTTP API, \texttt{TransferError::InsufficientFunds} maps to HTTP \texttt{422 Unprocessable Entity} with a JSON body telling the client their available balance, and \textbf{must not} wake the on-call SRE at 3:00~AM.
  \item \texttt{TransferError::AccountFrozen} maps to HTTP \texttt{403 Forbidden} and emits a compliance audit event.
  \item \texttt{TransferError::Storage(sql\_err)} means PostgreSQL dropped a connection or timed out; it maps to HTTP \texttt{503 Service Unavailable}, hides internal SQL table names from external clients, and increments a critical Prometheus alert counter.
\end{itemize}

For Zone~1, define a concrete, exhaustive \texttt{enum} error per module or subsystem. Let us first write out the pure standard-library implementation so you see every trait involved, and then show the concise \texttt{\#[derive(thiserror::Error)]} macro used in almost every production Rust workspace:

\begin{lstlisting}[language=Rust]
use std::fmt;

#[derive(Debug)]
pub struct SqlError {
    pub code: &'static str,
    pub detail: String,
}

impl fmt::Display for SqlError {
    fn fmt(&self, f: &mut fmt::Formatter<'_>) -> fmt::Result {
        write!(f, "database error ({}): {}", self.code, self.detail)
    }
}
impl std::error::Error for SqlError {}

// Structured domain error enum for the transfer subsystem:
#[derive(Debug)]
pub enum TransferError {
    InsufficientFunds { available_cents: i64, requested_cents: i64 },
    AccountFrozen { account_id: u64 },
    SameSourceAndDestination { account_id: u64 },
    Storage(SqlError),
}

// 1. Human-readable single-level description (no trailing newline!)
impl fmt::Display for TransferError {
    fn fmt(&self, f: &mut fmt::Formatter<'_>) -> fmt::Result {
        match self {
            Self::InsufficientFunds {
                available_cents,
                requested_cents,
            } => write!(
                f,
                "insufficient funds: have {available_cents} cents, \
                 need {requested_cents} cents"
            ),
            Self::AccountFrozen { account_id } => {
                write!(f, "account {account_id} is frozen by compliance")
            }
            Self::SameSourceAndDestination { account_id } => {
                write!(f, "self-transfer forbidden on {account_id}")
            }
            Self::Storage(_) => write!(f, "ledger storage query failed"),
        }
    }
}

// 2. Causal chain hook so loggers can walk `.source()` down to SqlError!
impl std::error::Error for TransferError {
    fn source(&self) -> Option<&(dyn std::error::Error + 'static)> {
        match self {
            Self::Storage(err) => Some(err),
            _ => None,
        }
    }
}

// 3. Automatic `?` conversion from lower-level SqlError!
impl From<SqlError> for TransferError {
    fn from(err: SqlError) -> Self {
        Self::Storage(err)
    }
}
\end{lstlisting}

Notice how clean those three trait implementations are --- zero hidden runtime reflection, zero stack-trace capture overhead on the \texttt{InsufficientFunds} path. Because writing \texttt{Display}, \texttt{Error::source}, and \texttt{From} by hand across twenty enums gets repetitive, David Tolnay's \texttt{thiserror} crate generates \textbf{those exact three trait \texttt{impl} blocks at compile time} with zero runtime dependency overhead:

\begin{lstlisting}[language=Rust]
// Identical compiled machine code using `thiserror`:
#[derive(Debug, thiserror::Error)]
pub enum TransferError {
    #[error(
        "low funds: have {available_cents}, need {requested_cents}"
    )]
    InsufficientFunds { available_cents: i64, requested_cents: i64 },

    #[error("account {account_id} is frozen by compliance")]
    AccountFrozen { account_id: u64 },

    #[error("self-transfer forbidden on {account_id}")]
    SameSourceAndDestination { account_id: u64 },

    #[error("ledger storage query failed")]
    Storage(#[from] SqlError), // Generates From + Error::source()!
}
\end{lstlisting}

\begin{trapbox}
Do \textbf{not} slap \texttt{\#[from] std::io::Error} blindly onto a single variant if your function performs three \emph{different} I/O operations (say reading a fee config file, reading a TLS certificate, and writing an audit socket). A bare \texttt{io::Error("No such file or directory")} never records \emph{which} path failed. Instead, wrap underlying errors with operational fields:
\begin{lstlisting}[language=Rust]
#[error("failed reading fee schedule at {path}")]
FeeConfigRead {
    path: std::path::PathBuf,
    #[source]
    source: std::io::Error,
}
\end{lstlisting}
and map at the call site with \texttt{.map\_err(|source| ConfigError::FeeConfigRead \{ path: p.to\_path\_buf(), source \})?}.
\end{trapbox}

\subsection{Zone 2: Application Entrypoints, Workers, and CLIs Use \texttt{anyhow}}

At the top of a binary crate --- \texttt{fn main()}, a nightly settlement batch job, a database migration CLI, or an integration test harness --- nobody above you is going to \texttt{match} on variants. The job of the top level is to glue together five subsystems (CLI flags, TOML files, TLS certs, PostgreSQL pools, HTTP listeners) and, if startup or batch execution fails, emit a clear causal story for the operator on call.

That is what \texttt{anyhow::Error} (a compact single-word boxed trait-object error with causal context chaining and optional backtrace capture) and its \texttt{Context} extension trait are built for:

\begin{lstlisting}[language=Rust]
use anyhow::{Context, Result};
use std::path::Path;

pub fn bootstrap_fee_engine(config_path: &Path) -> Result<FeeSchedule> {
    let raw_toml = std::fs::read_to_string(config_path).with_context(|| {
        format!("failed reading fee config file {}", config_path.display())
    })?;

    let flat_cents: i64 = raw_toml
        .trim()
        .parse()
        .context("fee config must contain an integer cent amount")?;

    anyhow::ensure!(
        flat_cents >= 0,
        "flat wire fee cannot be negative (got {flat_cents})"
    );
    Ok(FeeSchedule { flat_wire_fee_cents: flat_cents })
}
\end{lstlisting}

Look at two subtle details in \texttt{bootstrap\_fee\_engine}:
\begin{enumerate}
  \item \textbf{Prefer \texttt{.with\_context(|| format!(...))} over \texttt{.context(format!(...))}.} Remember eager argument evaluation: \texttt{.context(format!(...))} formats and heap-allocates the error string \emph{even on the 99.99-percent success path}. Passing a closure to \texttt{.with\_context(|| ...)} allocates nothing unless an \texttt{Err} actually occurred. Use plain \texttt{.context("...")} only for zero-allocation \texttt{\&'static str} literals.
  \item \textbf{Causal chain output.} When \texttt{/etc/ledger/fees.toml} has bad permissions, formatting \texttt{err} with \texttt{\{err:\#\}} yields:
  \begin{quote}
  \small\texttt{failed reading fee config file /etc/ledger/fees.toml:}\\
  \texttt{Permission denied (os error 13)}
  \end{quote}
  capturing both the high-level business intent and the low-level OS kernel cause in a single log line.
\end{enumerate}

\subsection{The Hard Line Between \texttt{Result} and \texttt{panic!}}

Newcomers sometimes treat \texttt{.unwrap()} as ``quick error handling'' during prototyping and accidentally ship those calls to production. In a multi-threaded Tokio/Axum server, a \texttt{panic!} inside an HTTP handler unwinds and aborts that request task (turning a malformed user payload into an instant HTTP \texttt{500} crash or dropped socket, and potentially poisoning a \texttt{std::sync::Mutex} if a lock guard was held across the panic).

Keep the boundary sharp:

\begin{focusbox}
\begin{itemize}
  \item \textbf{Return \texttt{Result<T, E>}} whenever failure can be triggered by \emph{anything outside the current function's proven invariants}: network packets, HTTP bodies/headers, CLI args, environment variables, filesystem state, clock drift, database errors, or violated business rules (insufficient balance, expired quote).
  \item \textbf{Use \texttt{panic!}, \texttt{assert!}, or \texttt{.expect("invariant: ...")}} \emph{only} when failure means \textbf{a bug already exists in our own program} --- a broken internal invariant where continuing could corrupt state or miscalculate money.
\end{itemize}
\end{focusbox}

Notice why seasoned Rust engineers forbid bare \texttt{.unwrap()} in production code (often enforcing \texttt{\#[deny(clippy::unwrap\_used)]}) while permitting \texttt{.expect("invariant: ...")} in the rare spots where human reasoning exceeds what the type checker knows:
\begin{lstlisting}[language=Rust]
// Compile-time constant IP/header parsing or mathematically proven step:
let loopback: std::net::IpAddr = "127.0.0.1"
    .parse()
    .expect("invariant: literal '127.0.0.1' is valid IPv4 syntax");
\end{lstlisting}
Writing the word \texttt{"invariant: ..."} forces the author to articulate \emph{why} \texttt{None}/\texttt{Err} is impossible, and gives any future maintainer an immediate explanation if a refactor ever violates that assumption.

% ======================================================================
\section{Modules, Visibility Capsules, \texttt{\#[must\_use]}, and Rust Testing Culture}
\label{ch08:modules-and-testing}
% ======================================================================

All five patterns we have built so far --- arenas, newtypes, typestates, view structs, and structured errors --- depend on one foundational mechanism to keep their guarantees: \textbf{module visibility boundaries}.

\subsection{Privacy Lives at the Module Boundary, Not the Struct Boundary}

Here is a rule that catches every Java and C++ programmer off guard the first week they write Rust:

\begin{trapbox}
In Java and C++, \texttt{private} means \emph{private to the enclosing \texttt{class}}. Two classes in the same file cannot touch each other's \texttt{private} fields (unless C++ explicitly grants \texttt{friend}).

In Rust, \textbf{the unit of encapsulation is the \texttt{mod} (module), NOT the \texttt{struct}.} Every item and struct field in Rust is private to its enclosing module by default --- which means:
\begin{enumerate}
  \item Any function or sibling struct written inside \emph{the same module} (and any \emph{descendant} child module, such as \texttt{mod tests} inside that file) \textbf{can read and mutate private fields directly}.
  \item Code in \emph{parent} or \emph{cousin} modules \textbf{cannot} access private items or construct structs with private fields unless the module explicitly marks them \texttt{pub}, \texttt{pub(crate)}, or \texttt{pub(super)}.
\end{enumerate}
\end{trapbox}

Why did Rust choose module-level privacy? Think back to \texttt{PostingView} or \texttt{WireOrder<S>}. Often two or three tightly coupled types (\texttt{Arena<T>}, \texttt{Slot<T>}, and \texttt{ArenaIter<'a, T>}) share a single representation invariant. Module-level privacy lets them cooperate without C++ \texttt{friend} declarations, \emph{provided you keep invariant-enforcing modules small}. If you dump 3,000 lines of unrelated business logic into one large \texttt{ledger.rs} file alongside \texttt{PositiveCents}, then all 3,000 lines can bypass \texttt{PositiveCents::try\_new}, wrap \emph{any} \texttt{NonZeroI64} directly (including a negative one), or mutate the private field.

Put core value types (\texttt{money.rs}, \texttt{iban.rs}, \texttt{wire\_state.rs}) in dedicated small leaf modules. Those modules act as \textbf{invariant capsules}: inside the 60-line module you audit every line that touches \texttt{.0}; across the remaining 50,000 lines of your workspace, the compiler enforces the invariant mechanically.

Alongside private (default) and \texttt{pub} (exported to the entire world), use \textbf{\texttt{pub(crate)}} liberally. Marking internal repository helpers, View Structs, or arena internals \texttt{pub(crate)} lets sibling modules inside your crate share them freely while keeping your crate's public API surface minimal and free to evolve without breaking external callers. Likewise, practice \emph{prelude restraint}: export explicit names rather than encouraging \texttt{use my\_crate::prelude::*;} wildcard imports that obscure where traits and types come from during code review.

\subsection{Compiler Guardrails with \texttt{\#[must\_use]}}

Have you ever wondered how \texttt{rustc} knows to print \texttt{warning: unused `Result` that must be used} whenever you forget \texttt{?} on a fallible call? There is no compiler special-casing for \texttt{Result}. In the standard library, \texttt{Result} is simply annotated with the attribute \texttt{\#[must\_use]}.

You can --- and should --- place \texttt{\#[must\_use = "reason"]} on your own pure value transformations, builders, and domain receipts. Consider an immutable money operation or settlement receipt:

\begin{lstlisting}[language=Rust]
#[must_use = "SettlementReceipt must be persisted to the outbox"]
pub struct SettlementReceipt {
    pub wire_id: u64,
    pub settled_cents: i64,
}

#[derive(Copy, Clone, Debug, PartialEq, Eq)]
pub struct Cents(pub i64);

impl Cents {
    #[must_use = "with_fee returns a new Cents without modifying self"]
    pub const fn with_fee(self, fee: Cents) -> Cents {
        Cents(self.0 + fee.0)
    }
}
\end{lstlisting}

If a teammate accustomed to mutable Java objects writes \texttt{total.with\_fee(wire\_fee);} as a standalone statement expecting \texttt{total} to mutate in place, the compiler immediately flags the bug with your exact custom explanation.

\subsection{Crate Layering for Production Services}

As a Rust backend grows beyond ten thousand lines (as our financial backend will in Part~V), organizing code into a multi-crate \textbf{Cargo workspace} reinforces clean architecture at the dependency-graph level because Cargo strictly forbids circular dependencies between crates:

\begin{enumerate}
  \item \textbf{\texttt{ledger-domain} (Pure Core Crate):} Holds Newtypes (\texttt{AccountId}, \texttt{PositiveCents}, \texttt{Iban}), Typestate protocols (\texttt{WireOrder<S>}), Arena/chart-of-accounts logic, fee-allocation math, and structured \texttt{thiserror} domain errors. Most importantly, \texttt{ledger-domain} has \textbf{zero network, Tokio, HTTP, or SQL dependencies}. It compiles in seconds, and all its tests run deterministically in memory.
  \item \textbf{\texttt{ledger-db} (Storage Adapter Crate):} Depends on \texttt{ledger-domain} and \texttt{sqlx}. Translates PostgreSQL rows into validated domain types (\texttt{AnyWireOrder}, \texttt{PositiveCents}) and maps database errors into domain storage errors.
  \item \textbf{\texttt{ledger-api} (Transport \& Binary Entrypoint Crate):} Depends on \texttt{ledger-domain}, \texttt{ledger-db}, \texttt{axum}, \texttt{tokio}, and \texttt{anyhow}. Parses untrusted JSON payloads at the perimeter into \texttt{ledger-domain} smart constructors, invokes application workflows, and maps \texttt{TransferError} variants to HTTP status codes.
\end{enumerate}

\begin{figure}[H]
\centering
\begin{tikzpicture}[
    font=\small,
    layer/.style={draw, rounded corners=5pt, very thick},
    arr/.style={-{Stealth[length=2.4mm,width=1.8mm]}, thick}
]
  % Outer shell: Adapters (Axum HTTP + SQLx Postgres)
  \draw[layer, draw=quotientorange!85!black, fill=quotientorange!5] (0,0) rectangle (13.8,5.15);
  \node[font=\bfseries\small, text=quotientorange!85!black, anchor=north west] at (0.3,4.95) {Outer Boundary Crates: \texttt{ledger-api} (Axum/Tokio) \& \texttt{ledger-db} (SQLx/Postgres)};
  \node[font=\scriptsize, text=gray!85!black, anchor=north west] at (0.3,4.52) {Untrusted JSON/SQL rows parsed once at edge; \texttt{anyhow} startup context; HTTP status mapping};

  % Middle capsule: pub(crate) Domain Services
  \draw[layer, draw=sourceblue!85!black, fill=sourceblue!5] (1.35,0.55) rectangle (12.45,3.95);
  \node[font=\bfseries\small, text=sourceblue!85!black, anchor=north west] at (1.65,3.75) {\texttt{ledger-domain} Crate (\texttt{pub} APIs + \texttt{pub(crate)} View Structs \& Workflows)};
  \node[font=\scriptsize, text=gray!85!black, anchor=north west] at (1.65,3.32) {Zero I/O dependencies! Pure deterministic functions, \texttt{Arena<AccountNode>}, \texttt{thiserror} enums};

  % Inner invariant capsules: Small private leaf modules
  \draw[layer, draw=targetgreen!80!black, fill=white] (2.15,0.95) rectangle (6.85,2.75);
  \node[font=\bfseries\scriptsize, text=targetgreen!75!black] at (4.5,2.42) {Leaf Capsule: \texttt{mod money} / \texttt{mod iban}};
  \node[font=\scriptsize\ttfamily, align=center] at (4.5,1.65) {
    private .0 fields!\\
    PositiveCents::try\_new()\\
    Iban::from\_str()
  };

  \draw[layer, draw=structurepurple!85!black, fill=white] (7.25,0.95) rectangle (11.75,2.75);
  \node[font=\bfseries\scriptsize, text=structurepurple!85!black] at (9.5,2.42) {Leaf Capsule: \texttt{mod wire\_state}};
  \node[font=\scriptsize\ttfamily, align=center] at (9.5,1.65) {
    WireOrder<Draft>\\
    approve(self) -> Cleared\\
    settle(self) -> Settled
  };

  % Dependency arrows pointing strictly inward!
  \draw[arr, quotientorange!85!black] (0.65,2.25) -- (1.35,2.25);
  \draw[arr, quotientorange!85!black] (13.15,2.25) -- (12.45,2.25);
\end{tikzpicture}
\caption{Concentric invariant-capsule architecture. Untrusted inputs are parsed at the outer HTTP and database boundary into proof-carrying Newtypes and Typestates owned by small leaf modules inside a zero-I/O \texttt{ledger-domain} core.}
\label{ch08:fig-invariant-capsule}
\end{figure}

\subsection{Rust's Four-Tier Testing Culture and Conservation Laws}

Rust ships a first-class test runner inside \texttt{cargo test} with zero external framework setup, and idiomatic crates combine four complementary layers of verification:

\begin{enumerate}
  \item \textbf{Co-located Unit Tests (\texttt{\#[cfg(test)] mod tests}):} Placed at the bottom of each source file. The \texttt{\#[cfg(test)]} conditional-compilation gate completely strips the module from release builds, while its position as a child module allows testing private helper invariants alongside public constructors.
  \item \textbf{Black-Box Integration Tests (\texttt{tests/*.rs}):} Every file placed in a crate's top-level \texttt{tests/} directory is compiled by Cargo as a \emph{separate external crate} linking against your library. It cannot touch private or \texttt{pub(crate)} items --- proving that external callers can actually use your public API cleanly and cannot subvert capsule invariants.
  \item \textbf{Executable Documentation Tests (\texttt{///} doc-comments):} Every fenced Rust snippet inside a \texttt{///} doc-comment is compiled and executed during \texttt{cargo test --doc}. Documentation examples in Rust never silently drift out of date when signatures evolve.
  \item \textbf{Property and Conservation-Law Tests:} Financial math has sharp algebraic invariants (such as \emph{conservation of money}: splitting a fee across $N$ sub-accounts must never create or destroy a single cent).
\end{enumerate}

Let us see conservation-law testing in action on a classic ledger problem: splitting a total fee of \texttt{total\_cents} equally across \texttt{n} participating sub-accounts without floating-point rounding drift. Using integer division and remainder (\texttt{div\_euclid} and \texttt{rem\_euclid}), the first \texttt{rem} accounts receive \texttt{q + 1} cents and the remaining \texttt{n - rem} accounts receive \texttt{q} cents:

\begin{lstlisting}[language=Rust]
/// Splits `total_cents >= 0` across `buckets > 0` accounts so the shares
/// sum to `total_cents` to the exact cent with at most 1-cent difference.
///
/// ```
/// # use my_ledger::allocate_fee_shares;
/// assert_eq!(allocate_fee_shares(100, 3), vec![34, 33, 33]);
/// ```
pub fn allocate_fee_shares(total_cents: i64, buckets: usize) -> Vec<i64> {
    assert!(total_cents >= 0, "invariant: fee total must be non-negative");
    assert!(buckets > 0, "invariant: bucket count must be positive");

    let n = buckets as i64;
    let base = total_cents.div_euclid(n);
    let remainder = total_cents.rem_euclid(n) as usize;

    (0..buckets)
        .map(|i| if i < remainder { base + 1 } else { base })
        .collect()
}

#[cfg(test)]
mod tests {
    use super::*;

    #[test]
    fn fee_allocation_conserves_exact_cents_across_sweep() {
        // Verify conservation + fairness across 5,000 inputs
        // (For huge input spaces, pair this with `proptest`!)
        for total in 0..=250_i64 {
            for buckets in 1..=20_usize {
                let shares = allocate_fee_shares(total, buckets);
                assert_eq!(shares.len(), buckets);

                // Law 1: Exact cent conservation (0 cents gained/lost)
                let sum: i64 = shares.iter().copied().sum();
                assert_eq!(sum, total);

                // Law 2: Fairness (max share - min share <= 1 cent)
                let min = *shares.iter().min().expect("buckets >= 1");
                let max = *shares.iter().max().expect("buckets >= 1");
                assert!(max - min <= 1);
            }
        }
    }
}
\end{lstlisting}

\subsection{Summary: Translating Instincts into Idiomatic Rust}

Table~\ref{ch08:tab-antipatterns} condenses every architectural shift from this chapter into a quick desk reference you can consult whenever code starts feeling awkward or borrow errors pile up.

\begin{table}[H]
\centering
\small
\begin{tabular}{@{}p{4.3cm}p{4.4cm}p{4.9cm}@{}}
\toprule
\textbf{Java~6 / C++ Instinct} & \textbf{Idiomatic Rust Pattern} & \textbf{Hardware \& Type Payoff} \\
\midrule
Cyclic \texttt{Rc<RefCell<Node>>} object graphs & Contiguous \texttt{Vec} Arena + \texttt{Copy} generational keys & Sequential cache lines, zero cycle leaks, no \texttt{BorrowMutError} \\
\addlinespace[3pt]
Raw \texttt{u64}/\texttt{i64}/\texttt{\&str} params + repeated \texttt{if} checks & \texttt{\#[repr(transparent)]} Newtypes + smart constructors & Register-identical ABI, 0-cost compile-time proof tokens \\
\addlinespace[3pt]
Mutable \texttt{boolean} flags + \texttt{IllegalStateException} & \texttt{Order<S>} Typestate + consuming \texttt{self} moves & Illegal order \& double-execution caught before binary runs \\
\addlinespace[3pt]
Heterogeneous DB rows fighting Typestate & \texttt{enum AnyOrder} wrapping \texttt{Order<S>} variants & Runtime DB hydration + static safety inside \texttt{match} arms \\
\addlinespace[3pt]
\texttt{\&String}, \texttt{\&Vec<T>}, \texttt{\&Box<T>} function params & \texttt{\&str}, \texttt{\&[T]}, \texttt{\&T}, \texttt{impl AsRef<Path>} & Removes double pointer hop; accepts literals \& sub-slices \\
\addlinespace[3pt]
Always \texttt{.to\_string()} on mostly clean strings & \texttt{Cow<'a, str>} (\texttt{Borrowed} vs \texttt{Owned}) & Zero \texttt{malloc} on 99-percent clean fast path \\
\addlinespace[3pt]
God-struct \texttt{\&mut self} locking disjoint fields & Sub-struct methods or explicit View Structs & Modular field-level borrowing without \texttt{RefCell} \\
\addlinespace[3pt]
Stringly errors or \texttt{.unwrap()} on user input & \texttt{thiserror} in domain, \texttt{anyhow} at app top & Exhaustive HTTP error mapping + rich operator causal logs \\
\bottomrule
\end{tabular}
\caption{Summary of anti-patterns and their idiomatic Rust architectural replacements.}
\label{ch08:tab-antipatterns}
\end{table}

Part~V uses each of these patterns directly: generational handles for in-memory account graphs, validated newtypes for amounts and IBANs at the HTTP boundary, typestate for transfer settlement, \texttt{thiserror} and \texttt{anyhow} for the two error zones, and conservation-law tests for the ledger.

% ======================================================================
\section{Exercises}
\label{ch08:exercises}
% ======================================================================

\begin{exercise}[Refactoring \texttt{Rc<RefCell<OrderNode>>} to a Generational Arena]
A teammate ported a Java~6 parent-child order-slicing graph (where cancelling a parent institutional order cancels all active child slice orders, and filling a child slice updates the parent's \texttt{remaining\_qty}) using \texttt{Rc<RefCell<OrderNode>>} and \texttt{Weak<RefCell<OrderNode>>}. Calling \texttt{fill\_slice} mutably borrows a child, upgrades \texttt{parent}, and calls \texttt{parent.\allowbreak{}borrow\_mut().\allowbreak{}reconcile\_children()}, triggering a runtime \texttt{BorrowMutError} panic.
\begin{enumerate}[label=(\alph*)]
  \item Replace the pointer graph with an \texttt{OrderArena} owning \texttt{Vec<Slot<OrderNode>>} and an 8-byte \texttt{Copy} handle \texttt{OrderKey \{ index: u32, generation: u32 \}}.
  \item Implement \texttt{insert}, \texttt{remove} (with embedded free-list and saturating generation increment), \texttt{get}, \texttt{get\_mut}, and \texttt{get\_two\_mut} using \texttt{split\_at\_mut}.
  \item Implement \texttt{fill\_\allowbreak{}slice(\&mut self,\allowbreak\ slice: OrderKey,\allowbreak\ qty: u64)} and write two unit tests: one verifying parent \texttt{remaining\_\allowbreak{}qty} rollup without panicking, and a second proving that removing a slice and inserting a new order into the recycled slot causes the old \texttt{OrderKey} to return \texttt{None} (ABA protection).
\end{enumerate}
\end{exercise}

\begin{exercise}[Zero-Cost ABA \texttt{RoutingNumber} Newtype]
US Automated Clearing House (ACH) transfers identify banks via a 9-digit ABA routing transit number $d_1 d_2 d_3 d_4 d_5 d_6 d_7 d_8 d_9$. A string is a valid ABA routing number if and only if it consists of exactly 9 ASCII digits (\texttt{'0'..='9'}) and satisfies the weighted checksum identity
\[
  \bigl(3(d_1 + d_4 + d_7) + 7(d_2 + d_5 + d_8) + (d_3 + d_6 + d_9)\bigr) \bmod 10 \;=\; 0.
\]
\begin{enumerate}[label=(\alph*)]
  \item Inside a dedicated \texttt{pub mod routing}, define a transparent wrapper \texttt{Routing\allowbreak{}Number(\allowbreak{}[u8; 9])} marked \texttt{\#[repr(\allowbreak{}transparent)]} with a \emph{private} inner array and a structured \texttt{Routing\allowbreak{}Error} enum with three variants: \texttt{Wrong\allowbreak{}Length(usize)}, a \texttt{NonDigit} variant recording the offending position and character, and \texttt{Checksum\allowbreak{}Mismatch}.
  \item Implement \texttt{FromStr}, \texttt{TryFrom<\&str>}, \texttt{Display}, and \texttt{as\_str(\&self) -> \&str}. (Hint: once parse-time checks prove all 9 bytes lie in \texttt{b'0'..=b'9'}, calling \texttt{from\_utf8(\&self.0)} followed by \texttt{.expect(\allowbreak{}"invariant: 9 ASCII digits")} cannot fail) Use \texttt{.rem\_euclid(10)} for the checksum remainder.
  \item Write unit tests confirming \texttt{size\_of::<RoutingNumber>() == 9}, accepting a valid routing number (\texttt{"111000025"}), and rejecting a single-digit typo and an adjacent-digit transposition. (Weights $3$ and $7$ catch every single-digit error but miss swaps of two digits that differ by exactly $5$; explain why.)
\end{enumerate}
\end{exercise}

\begin{exercise}[Compile-Time \texttt{BatchSettlement<S>} Typestate + Runtime Envelope]
Model an end-of-day clearing batch that progresses strictly through three states: \texttt{Open} (accepts \texttt{push\_posting(\&mut self,\allowbreak\ AccountId,\allowbreak\ PositiveCents)}), \texttt{Sealed} (carries a 32-byte SHA-256 batch digest and frozen total cents; forbids adding more postings), and \texttt{Submitted} (carries a clearinghouse acknowledgement ID).
\begin{enumerate}[label=(\alph*)]
  \item Implement \texttt{Batch\allowbreak{}Settlement<S>} with consuming \texttt{self} transitions: \texttt{seal(self)} returning \texttt{Result<\allowbreak{}Batch\allowbreak{}Settlement<Sealed>,\allowbreak\ Empty\allowbreak{}Batch\allowbreak{}Error>}, and \texttt{submit(self,\allowbreak\ ack\_id: u64)} returning a pair \texttt{(Batch\allowbreak{}Settlement<Submitted>,\allowbreak\ Settlement\allowbreak{}Receipt)}. Mark \texttt{Settlement\allowbreak{}Receipt} with \texttt{\#[must\_use]}.
  \item Explain in two sentences why taking \texttt{self} by value in \texttt{submit} prevents double-submission in a way a Java~6 \texttt{SealedBatch.submit()} method cannot.
  \item Define a runtime \texttt{pub enum AnyBatch} wrapping all three typestates so an SQL loader can return \texttt{Vec<AnyBatch>} from mixed database rows and dispatch only \texttt{Open} batches to \texttt{seal()}.
\end{enumerate}
\end{exercise}

\begin{exercise}[Zero-Allocation Fast Path with \texttt{Cow<'\_, str>}]
Ledger account codes follow canonical format \texttt{"US-OPERATING-01"}: non-empty ASCII uppercase letters, digits, and single \texttt{'-'} separators, with no leading/trailing spaces and no underscores \texttt{'\_'}. Upstream CSV files usually arrive already canonical, yet occasionally contain lowercase letters (\texttt{"us-operating-01"}) or underscores (\texttt{"US\_OPERATING\_01"}).
\begin{enumerate}[label=(\alph*)]
  \item Write \texttt{normalize\_\allowbreak{}account\_\allowbreak{}code(raw: \&str)} returning \texttt{Result<\allowbreak{}Cow<'\_, str>,\allowbreak\ Code\allowbreak{}Error>} so that it trims ASCII whitespace, validates syntax, and returns \texttt{Ok(\allowbreak{}Cow::\allowbreak{}Borrowed(\allowbreak{}trimmed))} with \textbf{zero heap allocations} whenever \texttt{trimmed} is already canonical.
  \item Only when lowercase ASCII letters or \texttt{'\_'} appear, allocate a single \texttt{String} of exact capacity, uppercasing letters and replacing \texttt{'\_'} with \texttt{'-'}, returning \texttt{Ok(Cow::Owned(...))}.
  \item Write a unit test checking \texttt{matches!(res, Cow::Borrowed(\_))} and comparing raw pointers via \texttt{std::ptr::eq(\allowbreak{}res.as\_bytes().as\_ptr(),\allowbreak\ input.as\_bytes().as\_ptr())} to prove that clean inputs share the exact same memory address with zero copying.
\end{enumerate}
\end{exercise}

\begin{exercise}[Split Borrows and Two-Zone Error Architecture]
Consider a \texttt{Reconciliation\allowbreak{}Runner} struct owning three fields: two maps from \texttt{AccountId} to \texttt{i64} named \texttt{ledger\_\allowbreak{}balances} and \texttt{bank\_\allowbreak{}tolerances}, and a \texttt{Vec<\allowbreak{}Discrepancy>} named \texttt{discrepancies}.
\begin{enumerate}[label=(\alph*)]
  \item Show why \texttt{check\_\allowbreak{}row(\&mut self,\allowbreak\ acct: Account\allowbreak{}Id,\allowbreak\ cents: i64,\allowbreak\ max\_\allowbreak{}drift: \&i64)} fails to compile when called as \texttt{self.\allowbreak{}check\_\allowbreak{}row(id, cents, tol)} while iterating over \texttt{\&self.\allowbreak{}bank\_\allowbreak{}tolerances}. Refactor \texttt{Reconciliation\allowbreak{}Runner} using an explicit \texttt{Reconcile\allowbreak{}View<'a>} borrow-projection struct so the batch loop compiles cleanly with zero cloning and zero \texttt{RefCell}.
  \item Design a structured Zone~1 \texttt{Reconcile\allowbreak{}Error} enum (implementing \texttt{Display}, \texttt{Error}, and \texttt{From<\allowbreak{}CsvParseError>}) for the library crate, and write a Zone~2 CLI entrypoint \texttt{run\_cli(\allowbreak{}csv\_path: \&Path) ->\allowbreak\ anyhow::\allowbreak{}Result<()>} using \texttt{.with\_\allowbreak{}context(|| ...)} so file-path context is attached lazily only on failure.
\end{enumerate}
\end{exercise}
